\documentclass{article}
\usepackage{enumerate}
\usepackage{amssymb}
\usepackage{amsmath}
\usepackage{amsthm}
\usepackage{latexsym}
\usepackage[T1]{fontenc}
\usepackage[latin1]{inputenc}

% Journal headers

\usepackage{fancyhdr}
\pagestyle{fancy}

\let\titlepageAuthor\author
\renewcommand{\author}[1]{\newcommand{\headerAuthor}{#1}\titlepageAuthor{#1}} 
\let\titlepageTitle\title
\renewcommand{\title}[1]{\newcommand{\headerTitle}{#1}\titlepageTitle{#1}} 

\lhead{Modeling Expanding Universes: \leftmark} \chead{} \rhead{\thepage} 
\lfoot{\copyright Guilford College Journal of Undergraduate Mathematics} \cfoot{} \rfoot{ - at www.jiblm.org}
\renewcommand{\headrulewidth}{0.4pt}
\renewcommand{\footrulewidth}{0.4pt}


\begin{document}


\title{Modeling Expanding Universes \\ in Euclidean Spaces}
\author{Suzanne M. Lea \\ E. E. Posey \\ Jerry E. Vaughan}
\date{\vspace{3cm}
\emph{Journal of Undergraduate Mathematics}
  \\ Department of Mathematics
  \\ Guilford College
  \\ Greensboro, North Carolina 27410 
  }
\maketitle
\thispagestyle{empty}



\newpage

\setcounter{tocdepth}{3} 
\tableofcontents
\thispagestyle{empty}


\newpage

\section*{Preface}
\addcontentsline{toc}{section}{\numberline{}Preface}
\markboth{Preface}{Preface}

\pagenumbering{roman}

The purpose of this monograph is to introduce, at the undergraduate level, Euclidean $n$-space, $E^n$, and briefly its infinite-dimensional version, $l^2$, using topics in physics as motivation and background. The physics we use is especially concerned with $4$-dimensional models of spacetime, and we have discussed several $3$-dimensional manifolds in $E^4$, including $3$-spheres, $3$-tori, and $3$-hyperboloids. We hope the reader will gain enhanced intuition into the Euclidean and analytical geometry of such figures. The material in this monograph can be used to supplement courses which cover calculus in dimensions $2$ and $3$.

The topic of expanding universes (or more precisely, the inflating balloon model of an expanding universe) provides background for our introduction of $E^{n}$, and the topics of wormholes and parallel universes provide background for our introduction of $l^{2}$. We use only basic analytic geometry of Euclidean spaces and elementary differential calculus (albeit in higher dimensions). We include a number of mathematical exercises of varying difficulty. While it is not our purpose to go deeply into matters of physics (for an elementary introduction to the physics of expanding universes, see [21]), we include throughout brief discussions of the relevant physics, and give numerous references (both elementary and advanced) for further reading.

We wish to thank Paul Duvall for suggestions concerning the Lemma in {\S}5, and Creative Services, Learning Resources Center, UNCG for help with the figures.

\begin{flushright}
Suzanne M. Lea and Jerry E. Vaughan \\
Department of Mathematics \\
University of North Carolina at Greensboro \\
Greensboro, NC 27412
\vspace{0.5cm}

E. E. Posey \\
4311 Dogwood Drive \\
Greensboro, NC 27410
\vspace{0.5cm}

October 15, 1994
\end{flushright}


\newpage

\section{Brief history of the expanding universe model}
\markboth{1. Brief History of Expanding Universe Model}{1. Brief History of Expanding Universe Model}

\pagenumbering{arabic}

Isaac Newton believed that our universe was infinite and static (neither expanding nor contracting). This notion held sway for about 200 years, and gained almost the status of a religious belief. When Einstein analyzed gravity and developed his general theory of relativity, he found that his mathematical model predicted an expanding universe. Since this contradicted Newton's theory of a static universe, Einstein promptly added to his model a constant factor, the cosmological constant, and explained in his 1917 paper that the only purpose of this constant was to keep the universe static [5, p. 214].

$\ $

In the late 1920s, an American astronomer, Edwin Hubble, found that light from stars in far-distant galaxies has a spectral red-shift(see [14] and [15, p. 753-765]). The red shift is interpreted as a Doppler shift caused by motion of the source. Because it is a shift to lower frequencies, the source, i.e. the stars and their parent galaxies, must be moving away from the earth. The amount of the shift indicates how fast the sources are moving away. When Hubble correlated the velocities of recession with the distance of the sources, he found that the velocities varied directly as their distances:
\begin{equation}
v = Hr
\end{equation}
where $r$ is the distance of the source from the earth, $v$ is the recession velocity of the source (i.e., $v = \frac{dr}{dt}$), and $H$ is a number now called Hubble's constant (despite the name, the Hubble ``constant'' is not thought to be constant over time [16, p.709]. It follows Hubble's formula that the further away a source is from the Earth, the faster it is receding from the Earth.

$\ $

Hubble's discovery of the expanding universe shattered the static models. Einstein removed the cosmological constant from his equations of general relativity, calling it one of the biggest blunders [10, p.44]. Cosmologists developed theories to account for the expansion. The theory generally accepted today considers the universe to be finite and expanding. The expansion started at the time of the Big Bang, a cosmic explosion 10 to 20 billion years ago.

$\ $

Models of expanding universes used by cosmologists are concerned with other physical properties as well as expansion, and are rather complex. In addition, they are not well constrained by observational data [12]; exotic models such as hyper-toroidal universes cannot be excluded. Standard models of the physical universe use at least $4$ dimensions, and in fact models using more than $4$ dimensions have been considered in physics (e.g., [9] and [11]). For further discussion about the discovery of the expanding universe, see [25, p. 165-194], and [24]. For a non-technical discussion of whether recession velocities are quantized, see [7] (quantized recession velocities imply discontinuous spaces).


\section{The inflating balloon model}
\markboth{2. Inflating Balloon Model}{2. Inflating Balloon Model}

The inflating balloon model of the expanding universe is well-known [16, p. 719]. In this model, our ordinary $3$-dimensional universe (at a specific time) corresponds to the ($2$-dimensional) surface of a spherical balloon. We mark (perhaps with a felt-tip pen) various dots on the surface of the balloon, and we label two of the dots $p$ and $q$. We then inflate the balloon, and observe that the distance between the dots (measured on the surface of the balloon) increases with the increase in radius of the balloon. Strictly speaking, the dots labeled $p$ and $q$ at time $t$ are not the same dots that carry the labels $p$ and $q$ at a later time $u > t$ because they are located not only at different times, but also at different points in $3$-space. (The path swept out by $p$ as it moves during the time that the balloon is being inflated is called the \emph{worldline} of $p$.)

We assume that the center of the balloon is not moved in $3$-space during the inflation, and we observe that the angle formed by two points and the center of the sphere does not appear to change as the balloon is inflated (this is perhaps obvious if we disregard the time axis, and view the balloons as being stacked up concentrically in $3$-space). In {\S}5, using calculus, we give a proof that in any such ``cross-sectionally spherical'' model, the angle is preserved as in the balloon model, and moreover (under a reasonable assumption about differentiability) exhibits the ``Hubble property'' of expanding universes (these terms are defined in {\S}3, and {\S}5).


\section{Definitions}
\markboth{3. Definitions}{3. Definitions}

Let $n$ be a positive integer. \emph{Euclidean $n$-space}, denoted $E^{n}$, is defined as the set of all ordered $n$-tuples $(x_{1}, \cdots, x_{n})$ of real numbers. We follow a standard convention of using a superscript to denote the dimension of a manifold. We use the term ``manifold'' informally to refer to subsets of $E^{n}$ in which we are interested in this monograph. (Examples of $2$-manifolds in $E^{3}$ are $2$-spheres, hyperboloids, toruses, and planes.)

A point $y$ in $E^{n}$ is on the $x_{1}$-axis provided $y_{i} = 0$ for $2 \leq i \leq n$. We may assume that the $x_{1}$-axis is a time axis and denote real numbers on the $x_{1}$-axis by such variables as $t$ or $u$ as well as $x_{1}$. If we wished, we could also assume that units of length and units of time are equivalent. Identical length and time units result when the speed of light $c = 1$; for example, if distance is measured in light years and time in years, $c = 1 lt.yr./yr.$.

The $(n-1)$-sphere centered at $(y_{1}, \cdots , y_{n})$ with radius $r>0$, denoted $S^{n-1} ((y_{1}, \cdots, y_{n}), r)$, is the set of all points $x$ in $E^{n}$ which satisfy the equation
\begin{equation*}
(x_{1} - y_{1})^{2} + (x_{2} - y_{2})^{2} + \cdots + (x_{n} - y_{n})^{2} = r^{2}.
\end{equation*}

In this monograph, we are only interested in $n$-spheres with centers on the $x_{1}$-axis. The $(n-1)$-sphere with center at the point $(t, 0, \cdots, 0)$ on the $x_{1}$-axis is the set of all points $x$ in $E^{n}$ which satisfy the equation
\begin{equation*}
(x_{1} - t)^{2} + x_{2}^{2} + \cdots + x_{n}^{2} = r^{2}.
\end{equation*}

A hyperplane in $E^{n+2}$, perpendicular to the $x_{1}$-axis at the point $(t, 0, \cdots, 0)$, is the set of all points $x$ in $E^{n+2}$ such that $x_{1} = t$, and is denoted by $P_{t}^{n+l}$. A subset (or more technically, an $(n+1)$-dimensional manifold) $M^{n+1}$ of $E^{n+2}$ is said to be cross-sectionally spherical if each hyperplane perendicu1ar to the $x_{1}$-axis intersects $M^{n+1}$ in an $n$-sphere with center at $(t, 0, \cdots, 0)$, or in a single point, or in the empty set. For each $t \epsilon E^{1}$ such that this intersection is nonempty, we define $\gamma(t) \geq 0$ to be the radius of this cross-sectional $n$-sphere. Thus
\begin{equation*}
S^{n}( (t, 0, \cdots, 0), \gamma(t) ) = P_{t}^{n+1} \cap M^{n+1}.
\end{equation*}

The function $\gamma(t)$ is called the \emph{radius function of the manifold}. To simplify the notation for the cross-sectional $n$-spheres, we let
\begin{equation*}
S_{r} = S_{n}( (t, 0, \cdots, 0), \gamma(t) )
\end{equation*}
(since the radius of each cross-sectional $n$-sphere is determined by the manifold, we do not need to explicitly include the radius in our notation). Thus, a cross-sectionally spherical manifold $M^{n+1}$ can be viewed as the union of a family of $n$-spheres:
\begin{equation*}
M^{n+1} = \cup \{ S_{t}^{n} : t \epsilon E^{1} \text{ and } \gamma(t) \text{ is defined} \}.
\end{equation*}

In the context of expanding universes, $M^{n+1}$ is called a \emph{spacetime manifold}, $E^{n+2}$ is called the \emph{ambient space}, and the spherical cross-sections $S_{t}^{n}$ are called the \emph{universe at time $t$}. Figure 1 gives an illustration of these concepts, in which the ambient space is $E^{3}$, the spacetime manifold $M^{2}$ is a hyperboloid whose equation is $x_{1}^{2} + \epsilon ^{2} = x_{2}^{2} + x_{2}^{2}$ (see {\S}4), and the universe at time $t$, S$_{t}^{1}$, is a circle (depicted at several different times).

%fig1
\begin{center}
%TeXCAD Picture [meu-fig1.pic]. Options:
%\grade{\on}
%\emlines{\off}
%\epic{\off}
%\beziermacro{\on}
%\reduce{\on}
%\snapping{\off}
%\quality{8.00}
%\graddiff{0.01}
%\snapasp{1}
%\zoom{4.0000}
\unitlength 1mm % = 2.85pt
\linethickness{0.4pt}
\ifx\plotpoint\undefined\newsavebox{\plotpoint}\fi % GNUPLOT compatibility
\begin{picture}(64,45.75)(0,0)
\qbezier(2.5,10.5)(28.5,27.13)(56.5,10.25)
\qbezier(2.5,35)(28.5,18.38)(56.5,35.25)
\qbezier(28.75,26.5)(32.75,22.38)(28.75,18.75)
\qbezier(28.64,26.52)(24.31,22.27)(28.81,18.74)
\qbezier(37.12,27.4)(41.1,22.98)(37.65,17.85)
\qbezier(19.37,27.4)(23.35,22.98)(19.9,17.85)
\qbezier(44.9,29.7)(50.65,22.27)(45.43,15.56)
\qbezier(12.53,29.7)(18.27,22.27)(13.06,15.56)
%\dashline{1}(45,15.38)(43.75,17.38)
\multiput(44.93,15.3)(-.03125,.05){10}{\line(0,1){.05}}
\multiput(44.3,16.3)(-.03125,.05){10}{\line(0,1){.05}}
%\end
%\dashline{1}(12.63,15.38)(11.38,17.38)
\multiput(12.55,15.3)(-.03125,.05){10}{\line(0,1){.05}}
\multiput(11.93,16.3)(-.03125,.05){10}{\line(0,1){.05}}
%\end
%\dashline{1}(43.75,17.38)(42.5,19.63)
\multiput(43.68,17.3)(-.03125,.05625){10}{\line(0,1){.05625}}
\multiput(43.05,18.43)(-.03125,.05625){10}{\line(0,1){.05625}}
%\end
%\dashline{1}(11.38,17.38)(10.13,19.63)
\multiput(11.3,17.3)(-.03125,.05625){10}{\line(0,1){.05625}}
\multiput(10.68,18.43)(-.03125,.05625){10}{\line(0,1){.05625}}
%\end
%\dashline{1}(42.5,19.63)(42.13,21.38)
\multiput(42.43,19.55)(-.03125,.14583){6}{\line(0,1){.14583}}
%\end
%\dashline{1}(10.13,19.63)(9.75,21.38)
\multiput(10.05,19.55)(-.03125,.14583){6}{\line(0,1){.14583}}
%\end
%\dashline{1}(42.13,21.38)(42,23.5)
\put(42.05,21.3){\line(0,1){.708}}
\put(41.97,22.72){\line(0,1){.708}}
%\end
%\dashline{1}(9.75,21.38)(9.63,23.5)
\put(9.68,21.3){\line(0,1){.708}}
\put(9.6,22.72){\line(0,1){.708}}
%\end
%\dashline{1}(42,23.5)(42.63,25.75)
\multiput(41.93,23.43)(.02976,.10714){7}{\line(0,1){.10714}}
\multiput(42.35,24.93)(.02976,.10714){7}{\line(0,1){.10714}}
%\end
%\dashline{1}(9.63,23.5)(10.25,25.75)
\multiput(9.55,23.43)(.02976,.10714){7}{\line(0,1){.10714}}
\multiput(9.97,24.93)(.02976,.10714){7}{\line(0,1){.10714}}
%\end
%\dashline{1}(42.63,25.75)(43.88,27.75)
\multiput(42.55,25.68)(.03125,.05){10}{\line(0,1){.05}}
\multiput(43.18,26.68)(.03125,.05){10}{\line(0,1){.05}}
%\end
%\dashline{1}(10.25,25.75)(11.5,27.75)
\multiput(10.18,25.68)(.03125,.05){10}{\line(0,1){.05}}
\multiput(10.8,26.68)(.03125,.05){10}{\line(0,1){.05}}
%\end
%\dashline{1}(43.88,27.75)(45,29.38)
\multiput(43.8,27.68)(.03125,.04514){12}{\line(0,1){.04514}}
\multiput(44.55,28.76)(.03125,.04514){12}{\line(0,1){.04514}}
%\end
%\dashline{1}(11.5,27.75)(12.63,29.38)
\multiput(11.43,27.68)(.03125,.04514){12}{\line(0,1){.04514}}
\multiput(12.18,28.76)(.03125,.04514){12}{\line(0,1){.04514}}
%\end
%\dashline{1}(37,27.5)(37.13,27.5)
\put(36.93,27.43){\line(1,0){.063}}
%\end
%\dashline{1}(19.25,27.5)(19.38,27.5)
\put(19.18,27.43){\line(1,0){.063}}
%\end
%\dashline{1}(37.13,27.5)(36,26.13)
\multiput(37.05,27.43)(-.03125,-.03819){12}{\line(0,-1){.03819}}
\multiput(36.3,26.51)(-.03125,-.03819){12}{\line(0,-1){.03819}}
%\end
%\dashline{1}(19.38,27.5)(18.25,26.13)
\multiput(19.3,27.43)(-.03125,-.03819){12}{\line(0,-1){.03819}}
\multiput(18.55,26.51)(-.03125,-.03819){12}{\line(0,-1){.03819}}
%\end
%\dashline{1}(36,26.13)(35.13,25)
\multiput(35.93,26.05)(-.033654,-.043269){13}{\line(0,-1){.043269}}
%\end
%\dashline{1}(18.25,26.13)(17.38,25)
\multiput(18.18,26.05)(-.033654,-.043269){13}{\line(0,-1){.043269}}
%\end
%\dashline{1}(35.13,25)(34.88,23.25)
\multiput(35.05,24.93)(-.0313,-.2188){4}{\line(0,-1){.2188}}
%\end
%\dashline{1}(17.38,25)(17.13,23.25)
\multiput(17.3,24.93)(-.0313,-.2188){4}{\line(0,-1){.2188}}
%\end
%\dashline{1}(34.88,23.25)(35,21.63)
\put(34.8,23.18){\line(0,-1){.813}}
%\end
%\dashline{1}(17.13,23.25)(17.25,21.63)
\put(17.05,23.18){\line(0,-1){.813}}
%\end
%\dashline{1}(35,21.63)(35.63,20.13)
\multiput(34.93,21.55)(.03125,-.075){10}{\line(0,-1){.075}}
%\end
%\dashline{1}(17.25,21.63)(17.88,20.13)
\multiput(17.18,21.55)(.03125,-.075){10}{\line(0,-1){.075}}
%\end
%\dashline{1}(35.63,20.13)(36.88,18.63)
\multiput(35.55,20.05)(.032051,-.038462){13}{\line(0,-1){.038462}}
\multiput(36.39,19.05)(.032051,-.038462){13}{\line(0,-1){.038462}}
%\end
%\dashline{1}(17.88,20.13)(19.13,18.63)
\multiput(17.8,20.05)(.032051,-.038462){13}{\line(0,-1){.038462}}
\multiput(18.64,19.05)(.032051,-.038462){13}{\line(0,-1){.038462}}
%\end
%\dashline{1}(36.88,18.63)(37.75,17.88)
\multiput(36.8,18.55)(.03646,-.03125){12}{\line(1,0){.03646}}
%\end
%\dashline{1}(19.13,18.63)(20,17.88)
\multiput(19.05,18.55)(.03646,-.03125){12}{\line(1,0){.03646}}
%\end
\put(29,22.5){\vector(0,1){20}}
\put(29,22.5){\vector(1,0){31.75}}
%\vector(29,22.5)(15,8.5)
\put(15,8.5){\vector(-1,-1){.07}}\multiput(29,22.5)(-.03373494,-.03373494){415}{\line(0,-1){.03373494}}
%\end
\put(29,21){\makebox(0,0)[cc]{$0$}}
\put(30.5,28.5){\makebox(0,0)[cc]{$\epsilon$}}
\put(29,45.75){\makebox(0,0)[cc]{$x_3$}}
\put(64,22.75){\makebox(0,0)[cc]{$x_1$}}
\put(13.75,7){\makebox(0,0)[cc]{$x_2$}}
\put(28.5,4.5){\makebox(0,0)[cc]{Figure 1}}
\end{picture}
\end{center}

Our space-time manifolds are embedded in the ambient space, and from the point of view of mathematics, this is natural and convenient. We note, however, that the physical universe is not thought to be embedded in an ambient space, and therefore questions about what lies beyond the ``edge'' of the universe cannot even be asked about the physical universe because there is no ``space'' beyond its ``edge''.

A model of spacetime without an ambient space is given in Exercise 5.11.


\section{Models for $(n+1)$-dimensional hyperboloids, embedded in $E^{n+2}$}
\markboth{4. Models for $(n+1)$-dimen. Hyperboloids}{4. Models for $(n+1)$-dimen. Hyperboloids}

To gain intuition about $n$-dimensional Euclidean spaces, it is common to consider the special cases of one, two, and three dimensions, and there are many expositions that follow this pattern. Two that are particularly well-know are the book \emph{Flatland} [1], and the more recent book \emph{Sphereland} [4], which extends these analogies to include the topic of expanding universes. We follow this pattern, and discuss our model in low dimension before giving the general $n$-dimensional case. 

We begin with the case $n=1$. Let $\epsilon$ denote a positive real number. Let $H^{2}$ denote a (connected) $2$-dimensional hyperboloid (i.e., hyperboloid of one sheet; see Figure 1) in $E^{3}$, defined by the equation
\begin{equation*}
x_{1}^{2} + \epsilon = x_{2}^{2} + x_{3}^{2}.
\end{equation*}

This $\epsilon$ is the minimum radius of the universe in this model, and occurs at time $t = 0$. We want this radius to be so small that at time $t = 0$, the universe is in the state called the big bang. To achieve that we use a value for $\epsilon$ which is less than the radius of a universe at which the laws of physics break down. (According to current theory, we would take $\epsilon$ less than the Planck length, which is approximately $1.616 \times 10^{-33}$ centimeters [16, p. 12].) This small value of $\epsilon$ allows the laws of physics to start anew after the big bang, and it appears to be meaningless to ask what they were previous to the big bang.

Some further notation: an \emph{affine space} $P^{k}$ in $E^{n}$, $k < n$ is a copy of $E^{k}$ in $E^{n}$ that may or may not contain the origin. If $k = n-1$, the affine space is called a \emph{hyperplane}. The symbol $P^{k}$ usually has a subscript. In this monograph, hyperplanes are usually chosen perpendicular to an axis of $E^{n}$ or perpendicular to a circle with center at the origin.

Let $t$ be a fixed number on the $x_{1}$-axis. Let $P_{t}^{2}$ be the hyperplane containing the point $(t, 0, 0)$ and perpendicular to the $x_{1}$-axis.

\begin{description}
\item[Exercise 4.1:] Show that for each $t$ on the $x_{1}$-axis, $H^{2} \cap P_{t}^{2} = S_{t}^{1}$ is a circle with center at $(t, 0, 0)$ in $E^{3}$ and radius of length $\sqrt{t^{2} + \epsilon^{2} }$.
\end{description}

In this model, ``space'' at time $t$ is the $1$-sphere (circle) $S_{t}^{1}$, and the universe is contracting for $t < 0$, and expanding for $t > 0$. The minimum radius of $S_{t}^{1}$ is $\epsilon$, and the minimum occurs at $t = 0$. (Negative values of $t$ can be considered as designating time before this minimum radius occurred.) It is easy to visualize the circle $S_{t}^{1}$ expanding or contracting in $H^{2}$ embedded in $E^{3}$ (see Figure 1). Unfortunately, inhabitants of this model universe must live on the $1$-dimensional circumference of a circular disk in a manner similar to the creatures in \emph{Lineland} [1]. 

$\ $

To attempt more realism, we now consider the case $n=2$; i.e., we move up one dimension.

Let $H^{3}$ denote the $3$-dimensional hyperboloid in $E^4$ defined by the equation
\begin{equation*}
x_{1}^{2} + \epsilon^{2} = x_{2}^{2} + x_{3}^{2} + x_{4}^{2}.
\end{equation*}

Let $t$ be fixed, and let $P_{t}^{3}$ be the hyperplane containing $(t, 0, 0, 0)$ and perpendicular to the $x_{1}$-axis of $E^{4}$. In this case, ``space'' at any time $t$ is the ordinary sphere ($2$-sphere) $H^{3} \cap P_{t}^{3} = S_{t}^{2}$ with center at $(t, 0, 0, 0)$ and radius of length $\sqrt{t^{2} + \epsilon^{2}}$.

In order to visualize $S_{t}^{2}$ contracting in $H^{3}$ for $t<0$, and expanding in $H^{3}$ for $t>0$, it is helpful to recall that the center $(t, 0, 0, 0)$ of $S$ is on the $x_{1}$-axis which is perpendicular to the (space) axes that actually define $s_{1}$, and that $(t,0,0,0)$ is not on $H^{3}$. Inhabitants of this model, like those in \emph{Sphereland} [4], are confined to the surface of a $3$-ball (ordinary sphere) [4]. 

$\ $

We now move up one more dimension, to the case $n=3$.

Let $H^{4}$ denote a $4$-dimensional hyperboloid in $E^{5}$ defined by the equation
\begin{equation*}
x_{1}^{2} + \epsilon_{2}^{2} = x_{2}^{2} + x_{3}^{2} + x_{4}^{2} + x_{5}^{2}.
\end{equation*}

Let $t$ be fixed, and let $P_{t}^{4}$ be the hyperplane containing $(t, 0, 0, 0, 0)$ and perpendicular to the $x_{1}$-axis of $E^{5}$. In this case, ``space'' at any time $t$ is the sphere ($3$-sphere) $H^{4} \cap P_{t}^{3} = S_{t}^{3}$ with center at $(t, 0, 0, 0, 0)$ with radius of length $\sqrt{t^{2} + \epsilon^{2}}$. (note that the equation of the $3$-dimensional sphere $S_{t}^{3}$ involves 4 variables, and therefore is not easy to visualize directly). It is compatible with known physical observations that the universe in which we live could be a $3$-sphere expanding in a $4$-dimensional hyperboloid $H^{4}$ that is embedded in $5$-dimensional Euclidean space $E^{5}$. This $3$-sphere is our traditional $3$-dimensional space. Analogies in lower dimensions give us some intuition about this model. Indeed, there is no mathematical reason to stop this discussion with $E^{5}$, and in fact, physical models of the universe involving many more than $4$ dimensions have been considered (such as string theories; see [9], [11]).

$\ $

We now discuss the general case of dimension $n$.

Let $H^{n+1}$ denote an $(n+1)$-dimensional hyperboloid in $E^{n+2}$ defined by the equation
\begin{equation*}
x_{1}^{2} + \epsilon^{2} = x_{2}^{2} + x_{3}^{2} + \cdots + x_{n+1}^{2} + x_{n+2}^{2}.
\end{equation*}

Let $t$ be fixed, and let $P_{t}^{4}$ be the hyperplane containing $(t, 0, \cdots, 0)$ and perpendicular to the $x_{1}$-axis of $E^{n+2}$.

\begin{description}
\item[Exercise 4.2:] Show that $H^{n+1} \cap P_{t}^{n+1} = S_{t}^{n}$ is an $n$-sphere with center at $(t, 0, \cdots, 0)$ with radius $\gamma(t)$ of length $\sqrt{ t^{2} + \epsilon^{2} }$.
\end{description}



\section{Expanding Universes and the Hubble constant}
\markboth{5. Expanding Universes, Hubble Constant}{5. Expanding Universes, Hubble Constant}

We begin by establishing some notation. Let $M^{n+1}$ be a cross-sectionally spherical manifold in $E^{n+2}$ with radius function $\gamma(t)$. For any point $p \epsilon M^{n+1}$, there is a unique $t \epsilon E^{1}$ such that $p \epsilon S_{t}^{n}$, hence $p = p(t) = (t, p_{2}(t), \cdots, p_{n+2}(t) )$.

We now define the worldline of a point. Physically, a worldline gives the past, present, and future history of a point: where it is at a particular time for all times $t$. We assume the movement of a point through time is due exclusively to the expansion of the universe (cf {\S}2), and therefore given a point at two times $t_{0} < t_{1}$, it occupies the ``same position'' on the cross-sectional sphere at time $t_{1}$ as it does at time $t_{0}$. With this in mind, we define the worldline of a point $p \epsilon M^{n+1}$ with the aid of $S^{n}(O_{n+1},1)$, the unit $n$-sphere in $E^{n+1}$ centered at the origin $O_{n+1} \epsilon E^{n+1}$. For convenience of notation, we list the coordinates of a point $a \epsilon S^{n}(O_{n+1},1)$ as ``$2$ through $n+2$'' instead of ``$1$ through $n+1$''; i.e., $a = (a_{2}, \cdots, a_{n+2})$. We define the worldline determined by $a \epsilon S^{n}(O_{n+1},1)$ to be
\begin{equation*}
\Lambda(a) = \{ (t, \gamma(t) a_{2}, \cdots, \gamma(t) a_{n+2}) : {t \epsilon E}^{1} \}.
\end{equation*}

It is easy to see 
\begin{enumerate}[\hspace{1cm}(1)]
  \item $a \neq b$ if and only if $\Lambda(a) \cap \Lambda(b) = \emptyset$, and 
  \item for any $p \epsilon M^{n+1}$, $p \epsilon \Lambda(a)$ where $a = ( \frac{p_{2}(t)}{\gamma(t)}, \cdots, \frac{p_{n+2}(t)}{\gamma(t)} )$ (in other words, $\{ \Lambda(a) a \epsilon S^{n} (O_{n+1},1) \}$ is a partition of $M_{n+1})$. 
\end{enumerate}

We define the \emph{worldline of a point} $p \epsilon M_{n+1}$ to be $\Lambda(a)$ where 
\begin{equation*}
a = ( \frac{p_{2}(t)}{\gamma(t)}, \cdots, \frac{p_{n+2}(t)}{\gamma(t)} ).
\end{equation*}

This definition of worldline allows us to give an easy calculus proof of the lemma below, which says that the angle formed by two points $p(t), q(t)$ and the center $c(t)$ of the sphere $S_{t}^{n}$ does not change as the points move along their worldlines. In order to discuss this angle, let $p(t)$ and $q(t)$ be two points on $S_{t}^{n}$, and let $C_{t}$ be a great circle on $S_{t}^{n}$ through $p = p(t)$ and $q = q(t)$, and let $A_{pq}(t)$ be the shorter arc on $C_{t}$ with end points $p$ and $q$. We denote the length of this arc by $r(t)$; it is the distance between $p(t)$, $q(t)$ as measured on $S_{t}^{n}$. We define $\gamma(t)$ to be the angle of positive radian measure determined by $p$, $q$, and $c(t) = (t, 0, \cdots, 0)$ (with $c(t)$ as the vertex of the angle).

$\ $

The following Lemma is intuitively clear for an expanding circle (with fixed center). By choosing a great circle on an expanding $n$-sphere (with fixed center), it is again intuitively clear. The reader may still consider it intuitively clear when the center is moving as in the lemma. We give a formal proof.

\begin{description}
\item[Lemma:] If $M^{n+1}$ is a cross-sectionally spherical manifold with radius function $\gamma(t)$, and $\gamma(t) > 0$ for all $t$ in an open interval $J$, then the angle $\alpha(t)$ determined by $p(t)$, $q(t)$, and $c(t) = (t, 0, \cdots, 0)$ is a constant for all $t \epsilon J$.

\item[Proof:] Let $a,b$ $\epsilon S^{n}(O_{n+1}, 1)$ be such that $p(t) \epsilon \Lambda(a)$ and $q(t) \epsilon \Lambda(b)$. Using a standard formula from calculus, we show that $\alpha(t)$ is the angle between $a,b$ for all $t \epsilon J$. Put $v = p(t) - c(t)$ and $w = q(t) - c(t)$. Using a well-know formula from vector calculus, we have
  \begin{align*}
  \cos(\alpha(t)) & = \frac{v \cdot w}{\vert v \vert \ \vert w \vert} \\
                  & = \frac{ \gamma(t)^{2}\Sigma_{i=2}^{n+2}a_{i}b_{i} }{ \gamma(t)^2 } \\
                  & = a \cdot b
  \end{align*}
  which shows that $\gamma(t) = \angle aO_{n+1}b$ is independent of $t$.
\end{description}

\begin{description}
\item[Exercise 5.1:] Prove the Lemma by drawing a sketch and showing that appropriate triangles are similar (see Figure 1).
\end{description}

Our main definition is a version of Hubble's formula (1).

\begin{description}
\item[Definition 5.2:] A cross-sectionally spherical manifold
  \begin{equation}
  M^{n+1} = \cup \{ S_{n}^{t} : t \epsilon J \}
  \end{equation}
  is said to have the expanding Hubble property over the interval $J$ if there exists a real-valued function $H(t)$ such that $H(t) > 0$ for all $t \epsilon J$, and for any two points $p(t),q(t)$ in $S_{n}^{t}$, their distance apart $r(t)$ as measured on $S_{n}^{t}$ satisfies the relation $\frac{dr}{dt} = H(t) r(t)$.
\end{description}

We now discuss conditions on the radius function $\gamma(t)$ of a cross-sectionally spherical manifold which imply the Hubble property for the manifold.

\begin{description}
\item[Proposition:] Let $M^{n+1}$ be a cross-sectionally spherical manifold such that its radius function $\gamma(t) > 0$ for all $t$ in an open interval $J$. If $\frac{d\gamma}{dt} > 0$ over $J$, then $M_{n+1} = \cup \{ S_{t}^{n}: t \epsilon J \}$ has the expanding Hubble property over $J$.

\item[Proof:] As above, let $S_{t}^{n} = P_{t}^{n+1} \cap M^{n+1}$ for $t \epsilon J$ be the family of cross-sectional $n$-spheres with center $c(t)$ and radius $\gamma(t)$. Let $p(t), q(t)$ be distinct points on $S_{t}^{n}$.

  Let $\gamma(t)$ denote the angle of positive radian measure determined by $p(t)$, $q(t)$, and $(t, 0, \cdots, 0)$. By the Lemma, $\alpha(t) = \alpha$ (a constant) for all $t \epsilon J$. Let $r(t)$ denote the length of $A_{p(t)q(t)}$ (the distance between the two points as measured on $S_{t}^{n}$). Since $r(t)$ is the length of a sector of a circle, we have
  \begin{equation}
  r(t) = \alpha \gamma(t).
  \end{equation}

  Thus
  \begin{equation}
  \frac{dr}{dt} = \alpha \frac{d\gamma}{dt}.
  \end{equation}

  By substituting for $\alpha$, we get
  \begin{equation}
  \frac{dr}{dt} = \left( \frac{ d\gamma / dt }{ \gamma(t) } \right) r(t)
  \end{equation}
  or
  \begin{equation}
  \frac{dr}{dt} = H(t)r(t)
  \end{equation}

  thus $r(t)$ satisfies the desired relation where
  \begin{equation}
  H(t) = \left( \frac{ d\gamma / dt }{ \gamma(t) } \right)
  \end{equation}

  Clearly, $H(t)$ does not depend on $p(t)$ and $q(t)$. Since $\gamma(t)$ and $\frac{d\gamma}{dt}$ are positive, it follows that $H(t) > 0$, for all $t \epsilon J$, and this completes the proof.
\end{description}

Both the balloon analogy and our discussion in this section depend on the assumption that points $p$ and $q$ are receding from each other only because the universe is expanding. That is, we require $p$ and $q$ to be fixed on the surface of the balloon ($S_{t}^{n}$). In the case of the physical universe, this assumption can be justified under appropriate conditions [16, p. 709].

\begin{description}
\item[Exercise 5.3:] Give a discussion for collapsing universes analogous to the discussion for expanding universes. In particular, give the definition for ``contracting Hubble property'', and prove the result analogous to the above Proposition obtained by changing in the hypothesis of the above Proposition ``$\frac{d\gamma}{dt} > 0$'' to ``$\frac{d\gamma}{dt} < 0$'', and in the conclusion ``expanding'' to ``contracting''.
\end{description}

It follows from the Proposition that cross-sectionally spherical manifolds obey the version of Hubble's formula given in (6), and therefore they provide reasonable mathematical models of an expanding universe.

Since each radius function $\gamma(t)$ produces a model, we are reminded of a statement attributed to the English physicist J. J. Thomson [5, p. 242]: ``We have Einstein's space, de Sitter's space, expanding universes, contracting universes, vibrating universes, mysterious universes. In fact, the pure mathematician may create universes just by writing down an equation $\ldots$ he can have a universe of his own.''

\begin{description}

\item[Exercise 5.4:] Find $\gamma(t)$ and $H(t)$ in the spacetime manifold defined by equation (2). Sketch the manifold for several values of $n$, and describe what happens to the universe at time $t$ as $t$ varies from $0$ to $+\infty$, and from $-\infty$ to $+\infty$ (recall that we use the $x_{1}$-axis as the time axis).

\item[Exercise 5.5:] Show that in the following spacetime manifold $H(t) = H$ is indeed constant for all $t$ (cf. [7]):
  \begin{equation*}
  \epsilon^{2} \text{ exp}(2 Hx_{1}) = x_{2}^{2} + x_{3}^{2} + \cdots + x_{n+1}^{2} + x_{n+2}^{2}
  \end{equation*}
  where $H>0$ and $\epsilon > 0$ are constants and $\text{exp}(x)$ is the usual exponential function.

\item[Exercise 5.6:] Consider the spacetime manifold defined by
  \begin{equation*}
  1 = x_{2}^{2} + \cdots + x_{n+2}^{2}
  \end{equation*}

  Find $\gamma(t)$ and $H(t)$ in this manifold. Sketch the manifold for $n=1$, and describe what happens to the universe at time $t$ as $t$ varies. What geometric term describes the shape of this manifold?

\item[Exercise 5.7:] Consider the spacetime manifold defined by
  \begin{equation*}
  1 = x_{1}^{2} + \cdots + x_{n+2}^{2}
  \end{equation*}

  Find $\gamma(t)$ and $H(t)$ in this manifold. Sketch the manifold for $n=1$, and describe what happens to the universe at time $t$ as $t$ varies between $-1$ and $+1$. What can be said about the universe at time $t$ with $\vert t \vert > 1$? What geometric term describes the shape of this manifold?

\item[Exercise 5.8:] Define two manifolds in which the universe goes through a succession of expansions and collapses. In one manifold, have each collapse end in a point; and in the other manifold, have each collapse end in a sphere of arbitrarily small radius $\epsilon$ (see Figure 2).

%Fig2
\begin{center}
%TeXCAD Picture [meu-fig2.pic]. Options:
%\grade{\on}
%\emlines{\off}
%\epic{\off}
%\beziermacro{\on}
%\reduce{\on}
%\snapping{\off}
%\quality{8.00}
%\graddiff{0.01}
%\snapasp{1}
%\zoom{4.0000}
\unitlength 0.85mm % = 2.85pt
\linethickness{0.4pt}
\ifx\plotpoint\undefined\newsavebox{\plotpoint}\fi % GNUPLOT compatibility
\begin{picture}(123.9,45.06)(0,0)
\put(59.12,3.61){\makebox(0,0)[cc]{Figure 2}}
\qbezier(19.69,32.91)(24.45,27.34)(19.69,21.47)
\qbezier(61.76,32.91)(66.52,27.34)(61.76,21.47)
\qbezier(103.84,32.91)(108.59,27.34)(103.84,21.47)
\qbezier(26.23,20.28)(31.06,27.12)(26.08,33.95)
\qbezier(68.3,20.28)(73.13,27.12)(68.15,33.95)
\qbezier(110.38,20.28)(115.21,27.12)(110.23,33.95)
\qbezier(12.98,20.28)(17.81,27.12)(12.83,33.95)
\qbezier(55.05,20.28)(59.88,27.12)(54.9,33.95)
\qbezier(97.13,20.28)(101.96,27.12)(96.98,33.95)
\qbezier(32.62,17.01)(39.61,27.56)(33.51,38.12)
\qbezier(74.7,17.01)(81.68,27.56)(75.59,38.12)
\qbezier(116.77,17.01)(123.75,27.56)(117.66,38.12)
\qbezier(5.5,16.38)(12.48,26.94)(6.39,37.49)
\qbezier(47.57,16.38)(54.56,26.94)(48.46,37.49)
\qbezier(89.64,16.38)(96.63,26.94)(90.53,37.49)
\qbezier(.22,10.62)(18.2,32.62)(39.76,10.62)
\qbezier(42.29,10.62)(60.28,32.62)(81.83,10.62)
\qbezier(84.36,10.62)(102.35,32.62)(123.9,10.62)
\qbezier(.22,43.91)(18.2,21.91)(39.76,43.91)
\qbezier(42.29,43.91)(60.28,21.91)(81.83,43.91)
\qbezier(84.36,43.91)(102.35,21.91)(123.9,43.91)
%\dashline{1}(19.75,33)(18.63,31.63)
\multiput(19.68,32.93)(-.03111,-.03806){12}{\line(0,-1){.03806}}
\multiput(18.93,32.02)(-.03111,-.03806){12}{\line(0,-1){.03806}}
%\end
%\dashline{1}(61.82,33)(60.7,31.63)
\multiput(61.75,32.93)(-.03111,-.03806){12}{\line(0,-1){.03806}}
\multiput(61,32.02)(-.03111,-.03806){12}{\line(0,-1){.03806}}
%\end
%\dashline{1}(103.9,33)(102.77,31.63)
\multiput(103.83,32.93)(-.03139,-.03806){12}{\line(0,-1){.03806}}
\multiput(103.08,32.02)(-.03139,-.03806){12}{\line(0,-1){.03806}}
%\end
%\dashline{1}(18.63,31.63)(17.88,30.38)
\multiput(18.56,31.56)(-.03125,-.05208){12}{\line(0,-1){.05208}}
%\end
%\dashline{1}(60.7,31.63)(59.95,30.38)
\multiput(60.63,31.56)(-.03125,-.05208){12}{\line(0,-1){.05208}}
%\end
%\dashline{1}(102.77,31.63)(102.02,30.38)
\multiput(102.7,31.56)(-.03125,-.05208){12}{\line(0,-1){.05208}}
%\end
%\dashline{1}(17.88,30.38)(17.5,28.75)
\multiput(17.81,30.31)(-.03167,-.13583){6}{\line(0,-1){.13583}}
%\end
%\dashline{1}(59.95,30.38)(59.57,28.75)
\multiput(59.88,30.31)(-.03167,-.13583){6}{\line(0,-1){.13583}}
%\end
%\dashline{1}(102.02,30.38)(101.65,28.75)
\multiput(101.95,30.31)(-.03083,-.13583){6}{\line(0,-1){.13583}}
%\end
%\dashline{1}(17.5,28.75)(17.38,26.88)
\put(17.43,28.68){\line(0,-1){.935}}
%\end
%\dashline{1}(59.57,28.75)(59.45,26.88)
\put(59.5,28.68){\line(0,-1){.935}}
%\end
%\dashline{1}(101.65,28.75)(101.52,26.88)
\put(101.58,28.68){\line(0,-1){.935}}
%\end
%\dashline{1}(17.38,26.88)(17.88,24.75)
\multiput(17.31,26.81)(.03333,-.142){5}{\line(0,-1){.142}}
\multiput(17.64,25.39)(.03333,-.142){5}{\line(0,-1){.142}}
%\end
%\dashline{1}(59.45,26.88)(59.95,24.75)
\multiput(59.38,26.81)(.03333,-.142){5}{\line(0,-1){.142}}
\multiput(59.71,25.39)(.03333,-.142){5}{\line(0,-1){.142}}
%\end
%\dashline{1}(101.52,26.88)(102.02,24.75)
\multiput(101.45,26.81)(.03333,-.142){5}{\line(0,-1){.142}}
\multiput(101.78,25.39)(.03333,-.142){5}{\line(0,-1){.142}}
%\end
%\dashline{1}(17.88,24.75)(18.13,24)
\multiput(17.81,24.68)(.0312,-.0938){4}{\line(0,-1){.0938}}
%\end
%\dashline{1}(59.95,24.75)(60.2,24)
\multiput(59.88,24.68)(.0313,-.0938){4}{\line(0,-1){.0938}}
%\end
%\dashline{1}(102.02,24.75)(102.27,24)
\multiput(101.95,24.68)(.0313,-.0938){4}{\line(0,-1){.0938}}
%\end
%\dashline{1}(18.13,24)(19.88,21.5)
\multiput(18.06,23.93)(.033654,-.048077){13}{\line(0,-1){.048077}}
\multiput(18.93,22.68)(.033654,-.048077){13}{\line(0,-1){.048077}}
%\end
%\dashline{1}(60.2,24)(61.95,21.5)
\multiput(60.13,23.93)(.033654,-.048077){13}{\line(0,-1){.048077}}
\multiput(61,22.68)(.033654,-.048077){13}{\line(0,-1){.048077}}
%\end
%\dashline{1}(102.27,24)(104.02,21.5)
\multiput(102.2,23.93)(.033654,-.048077){13}{\line(0,-1){.048077}}
\multiput(103.07,22.68)(.033654,-.048077){13}{\line(0,-1){.048077}}
%\end
%\dashline{1}(25.88,33.88)(24.88,32.25)
\multiput(25.81,33.81)(-.03333,-.05433){10}{\line(0,-1){.05433}}
\multiput(25.14,32.72)(-.03333,-.05433){10}{\line(0,-1){.05433}}
%\end
%\dashline{1}(67.95,33.88)(66.95,32.25)
\multiput(67.88,33.81)(-.03333,-.05433){10}{\line(0,-1){.05433}}
\multiput(67.21,32.72)(-.03333,-.05433){10}{\line(0,-1){.05433}}
%\end
%\dashline{1}(110.02,33.88)(109.02,32.25)
\multiput(109.95,33.81)(-.03333,-.05433){10}{\line(0,-1){.05433}}
\multiput(109.28,32.72)(-.03333,-.05433){10}{\line(0,-1){.05433}}
%\end
%\dashline{1}(12.63,33.88)(11.63,32.25)
\multiput(12.56,33.81)(-.03333,-.05433){10}{\line(0,-1){.05433}}
\multiput(11.89,32.72)(-.03333,-.05433){10}{\line(0,-1){.05433}}
%\end
%\dashline{1}(54.7,33.88)(53.7,32.25)
\multiput(54.63,33.81)(-.03333,-.05433){10}{\line(0,-1){.05433}}
\multiput(53.96,32.72)(-.03333,-.05433){10}{\line(0,-1){.05433}}
%\end
%\dashline{1}(96.77,33.88)(95.77,32.25)
\multiput(96.7,33.81)(-.03333,-.05433){10}{\line(0,-1){.05433}}
\multiput(96.03,32.72)(-.03333,-.05433){10}{\line(0,-1){.05433}}
%\end
%\dashline{1}(24.88,32.25)(24.13,30.25)
\multiput(24.81,32.18)(-.03125,-.08333){8}{\line(0,-1){.08333}}
\multiput(24.31,30.85)(-.03125,-.08333){8}{\line(0,-1){.08333}}
%\end
%\dashline{1}(66.95,32.25)(66.2,30.25)
\multiput(66.88,32.18)(-.03125,-.08333){8}{\line(0,-1){.08333}}
\multiput(66.38,30.85)(-.03125,-.08333){8}{\line(0,-1){.08333}}
%\end
%\dashline{1}(109.02,32.25)(108.27,30.25)
\multiput(108.95,32.18)(-.03125,-.08333){8}{\line(0,-1){.08333}}
\multiput(108.45,30.85)(-.03125,-.08333){8}{\line(0,-1){.08333}}
%\end
%\dashline{1}(11.63,32.25)(10.88,30.25)
\multiput(11.56,32.18)(-.03125,-.08333){8}{\line(0,-1){.08333}}
\multiput(11.06,30.85)(-.03125,-.08333){8}{\line(0,-1){.08333}}
%\end
%\dashline{1}(53.7,32.25)(52.95,30.25)
\multiput(53.63,32.18)(-.03125,-.08333){8}{\line(0,-1){.08333}}
\multiput(53.13,30.85)(-.03125,-.08333){8}{\line(0,-1){.08333}}
%\end
%\dashline{1}(95.77,32.25)(95.02,30.25)
\multiput(95.7,32.18)(-.03125,-.08333){8}{\line(0,-1){.08333}}
\multiput(95.2,30.85)(-.03125,-.08333){8}{\line(0,-1){.08333}}
%\end
%\dashline{1}(24.13,30.25)(23.88,27.88)
\put(24.06,30.18){\line(0,-1){.79}}
\put(23.89,28.6){\line(0,-1){.79}}
%\end
%\dashline{1}(66.2,30.25)(65.95,27.88)
\put(66.13,30.18){\line(0,-1){.79}}
\put(65.96,28.6){\line(0,-1){.79}}
%\end
%\dashline{1}(108.27,30.25)(108.02,27.88)
\put(108.2,30.18){\line(0,-1){.79}}
\put(108.03,28.6){\line(0,-1){.79}}
%\end
%\dashline{1}(10.88,30.25)(10.63,27.88)
\put(10.81,30.18){\line(0,-1){.79}}
\put(10.64,28.6){\line(0,-1){.79}}
%\end
%\dashline{1}(52.95,30.25)(52.7,27.88)
\put(52.88,30.18){\line(0,-1){.79}}
\put(52.71,28.6){\line(0,-1){.79}}
%\end
%\dashline{1}(95.02,30.25)(94.77,27.88)
\put(94.95,30.18){\line(0,-1){.79}}
\put(94.78,28.6){\line(0,-1){.79}}
%\end
%\dashline{1}(23.88,27.88)(24,25.63)
\put(23.81,27.81){\line(0,-1){.75}}
\put(23.89,26.31){\line(0,-1){.75}}
%\end
%\dashline{1}(65.95,27.88)(66.07,25.63)
\put(65.88,27.81){\line(0,-1){.75}}
\put(65.96,26.31){\line(0,-1){.75}}
%\end
%\dashline{1}(108.02,27.88)(108.15,25.63)
\put(107.95,27.81){\line(0,-1){.75}}
\put(108.04,26.31){\line(0,-1){.75}}
%\end
%\dashline{1}(10.63,27.88)(10.75,25.63)
\put(10.56,27.81){\line(0,-1){.75}}
\put(10.64,26.31){\line(0,-1){.75}}
%\end
%\dashline{1}(52.7,27.88)(52.82,25.63)
\put(52.63,27.81){\line(0,-1){.75}}
\put(52.71,26.31){\line(0,-1){.75}}
%\end
%\dashline{1}(94.77,27.88)(94.9,25.63)
\put(94.7,27.81){\line(0,-1){.75}}
\put(94.79,26.31){\line(0,-1){.75}}
%\end
%\dashline{1}(24,25.63)(24.63,22.5)
\multiput(23.93,25.56)(.0315,-.1565){5}{\line(0,-1){.1565}}
\multiput(24.24,23.99)(.0315,-.1565){5}{\line(0,-1){.1565}}
%\end
%\dashline{1}(66.07,25.63)(66.7,22.5)
\multiput(66,25.56)(.0315,-.1565){5}{\line(0,-1){.1565}}
\multiput(66.31,23.99)(.0315,-.1565){5}{\line(0,-1){.1565}}
%\end
%\dashline{1}(108.15,25.63)(108.77,22.5)
\multiput(108.08,25.56)(.031,-.1565){5}{\line(0,-1){.1565}}
\multiput(108.39,23.99)(.031,-.1565){5}{\line(0,-1){.1565}}
%\end
%\dashline{1}(10.75,25.63)(11.38,22.5)
\multiput(10.68,25.56)(.0315,-.1565){5}{\line(0,-1){.1565}}
\multiput(10.99,23.99)(.0315,-.1565){5}{\line(0,-1){.1565}}
%\end
%\dashline{1}(52.82,25.63)(53.45,22.5)
\multiput(52.75,25.56)(.0315,-.1565){5}{\line(0,-1){.1565}}
\multiput(53.06,23.99)(.0315,-.1565){5}{\line(0,-1){.1565}}
%\end
%\dashline{1}(94.9,25.63)(95.52,22.5)
\multiput(94.83,25.56)(.031,-.1565){5}{\line(0,-1){.1565}}
\multiput(95.14,23.99)(.031,-.1565){5}{\line(0,-1){.1565}}
%\end
%\dashline{1}(24.63,22.5)(26,20.63)
\multiput(24.56,22.43)(.032619,-.044524){14}{\line(0,-1){.044524}}
\multiput(25.47,21.18)(.032619,-.044524){14}{\line(0,-1){.044524}}
%\end
%\dashline{1}(66.7,22.5)(68.07,20.63)
\multiput(66.63,22.43)(.032619,-.044524){14}{\line(0,-1){.044524}}
\multiput(67.54,21.18)(.032619,-.044524){14}{\line(0,-1){.044524}}
%\end
%\dashline{1}(108.77,22.5)(110.15,20.63)
\multiput(108.7,22.43)(.032857,-.044524){14}{\line(0,-1){.044524}}
\multiput(109.62,21.18)(.032857,-.044524){14}{\line(0,-1){.044524}}
%\end
%\dashline{1}(11.38,22.5)(12.75,20.63)
\multiput(11.31,22.43)(.032619,-.044524){14}{\line(0,-1){.044524}}
\multiput(12.22,21.18)(.032619,-.044524){14}{\line(0,-1){.044524}}
%\end
%\dashline{1}(53.45,22.5)(54.82,20.63)
\multiput(53.38,22.43)(.032619,-.044524){14}{\line(0,-1){.044524}}
\multiput(54.29,21.18)(.032619,-.044524){14}{\line(0,-1){.044524}}
%\end
%\dashline{1}(95.52,22.5)(96.9,20.63)
\multiput(95.45,22.43)(.032857,-.044524){14}{\line(0,-1){.044524}}
\multiput(96.37,21.18)(.032857,-.044524){14}{\line(0,-1){.044524}}
%\end
%\dashline{1}(33.38,38)(32,35.88)
\multiput(33.31,37.93)(-.03136,-.04818){11}{\line(0,-1){.04818}}
\multiput(32.62,36.87)(-.03136,-.04818){11}{\line(0,-1){.04818}}
%\end
%\dashline{1}(75.45,38)(74.07,35.88)
\multiput(75.38,37.93)(-.03136,-.04818){11}{\line(0,-1){.04818}}
\multiput(74.69,36.87)(-.03136,-.04818){11}{\line(0,-1){.04818}}
%\end
%\dashline{1}(117.52,38)(116.15,35.88)
\multiput(117.45,37.93)(-.03114,-.04818){11}{\line(0,-1){.04818}}
\multiput(116.76,36.87)(-.03114,-.04818){11}{\line(0,-1){.04818}}
%\end
%\dashline{1}(6.25,37.38)(4.88,35.25)
\multiput(6.18,37.31)(-.03114,-.04841){11}{\line(0,-1){.04841}}
\multiput(5.49,36.24)(-.03114,-.04841){11}{\line(0,-1){.04841}}
%\end
%\dashline{1}(48.32,37.38)(46.95,35.25)
\multiput(48.25,37.31)(-.03114,-.04841){11}{\line(0,-1){.04841}}
\multiput(47.56,36.24)(-.03114,-.04841){11}{\line(0,-1){.04841}}
%\end
%\dashline{1}(90.4,37.38)(89.02,35.25)
\multiput(90.33,37.31)(-.03136,-.04841){11}{\line(0,-1){.04841}}
\multiput(89.64,36.24)(-.03136,-.04841){11}{\line(0,-1){.04841}}
%\end
%\dashline{1}(32,35.88)(31.13,33.75)
\multiput(31.93,35.81)(-.03222,-.07889){9}{\line(0,-1){.07889}}
\multiput(31.35,34.39)(-.03222,-.07889){9}{\line(0,-1){.07889}}
%\end
%\dashline{1}(74.07,35.88)(73.2,33.75)
\multiput(74,35.81)(-.03222,-.07889){9}{\line(0,-1){.07889}}
\multiput(73.42,34.39)(-.03222,-.07889){9}{\line(0,-1){.07889}}
%\end
%\dashline{1}(116.15,35.88)(115.27,33.75)
\multiput(116.08,35.81)(-.03259,-.07889){9}{\line(0,-1){.07889}}
\multiput(115.49,34.39)(-.03259,-.07889){9}{\line(0,-1){.07889}}
%\end
%\dashline{1}(4.88,35.25)(4,33.13)
\multiput(4.81,35.18)(-.03259,-.07852){9}{\line(0,-1){.07852}}
\multiput(4.22,33.77)(-.03259,-.07852){9}{\line(0,-1){.07852}}
%\end
%\dashline{1}(46.95,35.25)(46.07,33.13)
\multiput(46.88,35.18)(-.03259,-.07852){9}{\line(0,-1){.07852}}
\multiput(46.29,33.77)(-.03259,-.07852){9}{\line(0,-1){.07852}}
%\end
%\dashline{1}(89.02,35.25)(88.15,33.13)
\multiput(88.95,35.18)(-.03222,-.07852){9}{\line(0,-1){.07852}}
\multiput(88.37,33.77)(-.03222,-.07852){9}{\line(0,-1){.07852}}
%\end
%\dashline{1}(31.13,33.75)(30.25,31.25)
\multiput(31.06,33.68)(-.03259,-.09259){9}{\line(0,-1){.09259}}
\multiput(30.47,32.01)(-.03259,-.09259){9}{\line(0,-1){.09259}}
%\end
%\dashline{1}(73.2,33.75)(72.32,31.25)
\multiput(73.13,33.68)(-.03259,-.09259){9}{\line(0,-1){.09259}}
\multiput(72.54,32.01)(-.03259,-.09259){9}{\line(0,-1){.09259}}
%\end
%\dashline{1}(115.27,33.75)(114.4,31.25)
\multiput(115.2,33.68)(-.03222,-.09259){9}{\line(0,-1){.09259}}
\multiput(114.62,32.01)(-.03222,-.09259){9}{\line(0,-1){.09259}}
%\end
%\dashline{1}(4,33.13)(3.13,30.63)
\multiput(3.93,33.06)(-.03222,-.09259){9}{\line(0,-1){.09259}}
\multiput(3.35,31.39)(-.03222,-.09259){9}{\line(0,-1){.09259}}
%\end
%\dashline{1}(46.07,33.13)(45.2,30.63)
\multiput(46,33.06)(-.03222,-.09259){9}{\line(0,-1){.09259}}
\multiput(45.42,31.39)(-.03222,-.09259){9}{\line(0,-1){.09259}}
%\end
%\dashline{1}(88.15,33.13)(87.27,30.63)
\multiput(88.08,33.06)(-.03259,-.09259){9}{\line(0,-1){.09259}}
\multiput(87.49,31.39)(-.03259,-.09259){9}{\line(0,-1){.09259}}
%\end
%\dashline{1}(30.25,31.25)(29.88,29.13)
\multiput(30.18,31.18)(-.0308,-.1767){4}{\line(0,-1){.1767}}
\multiput(29.93,29.77)(-.0308,-.1767){4}{\line(0,-1){.1767}}
%\end
%\dashline{1}(72.32,31.25)(71.95,29.13)
\multiput(72.25,31.18)(-.0308,-.1767){4}{\line(0,-1){.1767}}
\multiput(72,29.77)(-.0308,-.1767){4}{\line(0,-1){.1767}}
%\end
%\dashline{1}(114.4,31.25)(114.02,29.13)
\multiput(114.33,31.18)(-.0317,-.1767){4}{\line(0,-1){.1767}}
\multiput(114.08,29.77)(-.0317,-.1767){4}{\line(0,-1){.1767}}
%\end
%\dashline{1}(3.13,30.63)(2.75,28.5)
\multiput(3.06,30.56)(-.0317,-.1775){4}{\line(0,-1){.1775}}
\multiput(2.81,29.14)(-.0317,-.1775){4}{\line(0,-1){.1775}}
%\end
%\dashline{1}(45.2,30.63)(44.82,28.5)
\multiput(45.13,30.56)(-.0317,-.1775){4}{\line(0,-1){.1775}}
\multiput(44.88,29.14)(-.0317,-.1775){4}{\line(0,-1){.1775}}
%\end
%\dashline{1}(87.27,30.63)(86.9,28.5)
\multiput(87.2,30.56)(-.0308,-.1775){4}{\line(0,-1){.1775}}
\multiput(86.95,29.14)(-.0308,-.1775){4}{\line(0,-1){.1775}}
%\end
%\dashline{1}(29.88,29.13)(29.75,25.75)
\put(29.81,29.06){\line(0,-1){.845}}
\put(29.74,27.37){\line(0,-1){.845}}
%\end
%\dashline{1}(71.95,29.13)(71.82,25.75)
\put(71.88,29.06){\line(0,-1){.845}}
\put(71.81,27.37){\line(0,-1){.845}}
%\end
%\dashline{1}(114.02,29.13)(113.9,25.75)
\put(113.95,29.06){\line(0,-1){.845}}
\put(113.89,27.37){\line(0,-1){.845}}
%\end
%\dashline{1}(2.75,28.5)(2.63,25.13)
\put(2.68,28.43){\line(0,-1){.842}}
\put(2.62,26.74){\line(0,-1){.843}}
%\end
%\dashline{1}(44.82,28.5)(44.7,25.13)
\put(44.75,28.43){\line(0,-1){.842}}
\put(44.69,26.74){\line(0,-1){.843}}
%\end
%\dashline{1}(86.9,28.5)(86.77,25.13)
\put(86.83,28.43){\line(0,-1){.842}}
\put(86.76,26.74){\line(0,-1){.843}}
%\end
%\dashline{1}(29.75,25.75)(30.13,22.38)
\put(29.68,25.68){\line(0,-1){.843}}
\put(29.87,23.99){\line(0,-1){.843}}
%\end
%\dashline{1}(71.82,25.75)(72.2,22.38)
\put(71.75,25.68){\line(0,-1){.843}}
\put(71.94,23.99){\line(0,-1){.843}}
%\end
%\dashline{1}(113.9,25.75)(114.27,22.38)
\put(113.83,25.68){\line(0,-1){.843}}
\put(114.01,23.99){\line(0,-1){.843}}
%\end
%\dashline{1}(2.63,25.13)(3,21.75)
\put(2.56,25.06){\line(0,-1){.845}}
\put(2.74,23.37){\line(0,-1){.845}}
%\end
%\dashline{1}(44.7,25.13)(45.07,21.75)
\put(44.63,25.06){\line(0,-1){.845}}
\put(44.81,23.37){\line(0,-1){.845}}
%\end
%\dashline{1}(86.77,25.13)(87.15,21.75)
\put(86.7,25.06){\line(0,-1){.845}}
\put(86.89,23.37){\line(0,-1){.845}}
%\end
%\dashline{1}(30.13,22.38)(31.13,19.63)
\multiput(30.06,22.31)(.03333,-.09167){10}{\line(0,-1){.09167}}
\multiput(30.73,20.48)(.03333,-.09167){10}{\line(0,-1){.09167}}
%\end
%\dashline{1}(72.2,22.38)(73.2,19.63)
\multiput(72.13,22.31)(.03125,-.08594){8}{\line(0,-1){.08594}}
\multiput(72.63,20.93)(.03125,-.08594){8}{\line(0,-1){.08594}}
%\end
%\dashline{1}(114.27,22.38)(115.27,19.63)
\multiput(114.2,22.31)(.03125,-.08594){8}{\line(0,-1){.08594}}
\multiput(114.7,20.93)(.03125,-.08594){8}{\line(0,-1){.08594}}
%\end
%\dashline{1}(3,21.75)(4,19)
\multiput(2.93,21.68)(.03125,-.08594){8}{\line(0,-1){.08594}}
\multiput(3.43,20.3)(.03125,-.08594){8}{\line(0,-1){.08594}}
%\end
%\dashline{1}(45.07,21.75)(46.07,19)
\multiput(45,21.68)(.03125,-.08594){8}{\line(0,-1){.08594}}
\multiput(45.5,20.3)(.03125,-.08594){8}{\line(0,-1){.08594}}
%\end
%\dashline{1}(87.15,21.75)(88.15,19)
\multiput(87.08,21.68)(.03125,-.08594){8}{\line(0,-1){.08594}}
\multiput(87.58,20.3)(.03125,-.08594){8}{\line(0,-1){.08594}}
%\end
%\dashline{1}(31.13,19.63)(32.38,17.5)
\multiput(31.06,19.56)(.03125,-.05325){10}{\line(0,-1){.05325}}
\multiput(31.68,18.49)(.03125,-.05325){10}{\line(0,-1){.05325}}
%\end
%\dashline{1}(73.2,19.63)(74.45,17.5)
\multiput(73.13,19.56)(.03125,-.05325){10}{\line(0,-1){.05325}}
\multiput(73.75,18.49)(.03125,-.05325){10}{\line(0,-1){.05325}}
%\end
%\dashline{1}(115.27,19.63)(116.52,17.5)
\multiput(115.2,19.56)(.03125,-.05325){10}{\line(0,-1){.05325}}
\multiput(115.82,18.49)(.03125,-.05325){10}{\line(0,-1){.05325}}
%\end
%\dashline{1}(4,19)(5.25,16.88)
\multiput(3.93,18.93)(.03125,-.053){10}{\line(0,-1){.053}}
\multiput(4.55,17.87)(.03125,-.053){10}{\line(0,-1){.053}}
%\end
%\dashline{1}(46.07,19)(47.32,16.88)
\multiput(46,18.93)(.03125,-.053){10}{\line(0,-1){.053}}
\multiput(46.62,17.87)(.03125,-.053){10}{\line(0,-1){.053}}
%\end
%\dashline{1}(88.15,19)(89.4,16.88)
\multiput(88.08,18.93)(.03125,-.053){10}{\line(0,-1){.053}}
\multiput(88.7,17.87)(.03125,-.053){10}{\line(0,-1){.053}}
%\end
\qbezier(39.88,44)(41.13,45.06)(42.38,43.88)
\qbezier(81.95,44)(83.2,45.06)(84.45,43.88)
\qbezier(39.7,10.57)(40.95,9.51)(42.2,10.7)
\qbezier(81.77,10.57)(83.02,9.51)(84.27,10.7)
\end{picture}
\end{center}



\item[Exercise 5.9:] Prove that if $\frac{dr}{dt}(t) = H(t)r(t)$ where $H(t)$ is a Riemann integrable function, then there exists a function $F(t)$ and a constant $K$ such that $r(t) = K \text{ exp}(F(t))$.

\item[Exercise 5.10:] Let $r(t), s(t), H(t)$ be positive-valued functions of a real variable such that $\frac{dr}{dt} = H(t)r(t)$ and $\frac{ds}{dt} = H(t)s(t)$. Let $a_{1}, a_{2}$ be points in $E^{n}$ such that $\vert a_{1} \vert = \vert a_{2} \vert = 1$ (recall that $\vert x-y \vert$ denotes the Euclidean distance between points $x,y \epsilon E^{n}$; so we are saying that $a_{1}, a_{2}$ both have Euclidean distance $1$ from the origin). The two rays (half-lines) in $E^{n}$ emanating from the origin and passing through $a_{1}$, and $a_{2}$ can be parametrized by the functions $L_{1}(t) = r(t) \cdot a_{1}$ and $L_{2}(t) = s(t) \cdot a_{2}$. Prove that the function $R(t) = \vert L_{1}(t) - L_{2}(t) \vert$ satisfies the property $\frac{dR}{dt} = H(t)R(t)$.

\item[Exercise 5.11:] Construct models of a universe having the expanding Hubble property with spacetime $E^{4}$ and with the universe at time $t$, the $3$-ball $B_{t}^{3}$ such that if $u<t$, then $B_{u}^{3} \subset B_{t}^{3}$ (this model of spacetime does not have an ambient space). In this exercise, define worldlines to be rays emanating from the origin, measure the distance between two points $p(t), q(t)$ using the usual Euclidean distance (see Exercise 5.10), and treat the $x_4$-axis as the time axis. Hint: Start with two positive real-valued functions $s(t)$ and $H(t)$ such that $s(t)$ is unbounded, $\frac{ds}{dt} = H(t) s(t)$, and $\frac{ds}{dt} > 0$. Define $B_{t}^{3} = \cup \{ s(u)y : y \epsilon U^{2} \text{ and } 0 \leq u \leq t \}$, where $U^{2}$ is the unit $2$-sphere with center at the origin of $E^{3}$. First show that there exist points $a = (a_{1}, a_{2}, a_{3}), b = (b_{1}, b_{2}, b_{3}) \epsilon U^{2}$ and constants $0 \leq k_{1}, k_{2} \leq 1$ such that $p(t) = k_{1} s(t)a$ and $q(t) = k_{2} s(t)b$, and then consider the equation $R(t) = \vert p(t)-q(t) \vert$ (compare with Exercise 5.10).

\end{description}

The following exercises deal with manifolds which are toroidal in some sense. The most familiar example of a toroidal manifold is the $2$-torus (the surface of a doughnut) in $E^{3}$ (see [3, 1.1-1.3]).

An \emph{$n$-torus} $T^{n}$ is an $n$-dimensional manifold that is cross-sectionally spherical with respect to a circle $A$, called the axial circle. We consider only the case where the circle $A$ lies in the $x_{1}x_{2}$-plane, with center at the origin.

\begin{description}

\item[Exercise 5.12:] Let $T^{n+1}$ be the manifold defined by the equation
  \begin{equation*}
  [(x_{1}^{2} + x_{2}^{2})^{\frac{1}{2}} - R]^{2} + x_{3}^{2} + \cdots + x_{n+2}^{2} = r(t)^{2}
  \end{equation*}
  where $0 \leq r(t) < R$, and $R>0$ is a constant. Show that $T^{n+1}$ is an $n+1$-torus. Hint: for $0 \leq t \leq 2\pi$, define a family of hyper-half-planes by

  $P_{t}^{n+1} =$
  \begin{equation*}
  \begin{cases}
    \{ (x_{1}, x_{1} \tan(t), x_{3}, \cdots, x_{n+2}): x_{1} \geq 0 \} \text{  if } 0 \leq t \leq \frac{\pi}{2}, \text{ or } \frac{3\pi}{2} \leq t < 2\pi \\
    \{ (x_{1}, x_{1} \tan(t), x_{3}, \cdots, x_{n+2}): x_{1} \leq 0 \} \text{  if } \frac{\pi}{2} < t < \frac{3\pi}{2}
  \end{cases}
  \end{equation*}
  and show that the circle $A$, defined in the $x_{1}x_{2}$-plane with center at the origin and radius $R$, is the axial circle of $T^{n+1}$ by showing that $P_{t}^{n+1} \cap T^{n+1}$ is an $n$-sphere with center $p(t) = (a(t), b(t), 0, \cdots, 0)$ on $A$ and radius $r(t)$. (To consider $T^{n+1}$ as a model of spacetime, the time variable is $t$, which can be viewed as the angle formed by the line through the origin and the $x_{1}$-axis).

\item[Exercise 5.13:] Let $\hat{T}^{2}$ be the version of 5.12 with $n=1$, $R=4$, and

  \begin{equation*}
  r(t) =
  \begin{cases}
    \sin^{2} \frac{t}{2} + \epsilon \text{  if } 0 \leq t \leq 2 \pi \\
    \frac{(t - 2\pi)}{t} + \epsilon \text{  if } t > 2\pi
  \end{cases}
  \end{equation*}

  Sketch the manifold, and describe what happens to the universe at times $t=0$ and $t=2\pi$. Over what interval is the universe expanding (respectively, collapsing)?

\item[Exercise 5.14:] Let $\hat{T}^{2}$ be the version of 5.12 with $n=1$, $R=4$, and
  \begin{equation*}
  r(t) = 
  \begin{cases}
    2\sin^{2} \frac{t}{8} + \epsilon \text{  if } 0 \leq t \leq 8\pi \\
    \tan(t) + \epsilon \text{  if } 8\pi < t \leq 8\pi + \frac{\pi}{4} \\
    1 + \epsilon \text{  if } t > 8\pi + \frac{7\pi}{4}
  \end{cases}
  \end{equation*}

  Sketch the manifold, and describe the universe over the four time intervals $0 \leq t \leq 2\pi$, $2\pi \leq t \leq 4\pi$, $4\pi \leq t \leq 6\pi$, and $6\pi \leq t \leq 8\pi$. Hint: the part of the manifold defined over $0 \leq t \leq 2\pi$ is in the ``interior'' of the part of the manifold defined over $2\pi \leq t \leq 4\pi$.

\end{description}

In the following exercises, we introduce a cross-sectionally toroidal spacetime manifold, denoted $V^{4}$. For a constant $0 \leq k \leq 1$ and real-valued function $R(t)$ with $R(t) > 0$, $\frac{dR}{dt} > 0$, for all $t \geq 0$, we define
\begin{equation*}
T_{t}^{3} = \{ (x_{1}, x_{2}, x_{3}, x_{4}): [(x_{1}^{2} + x_{1}^{2})^{\frac{1}{2} } - R(t)]^{2} + x_{3}^{2} + x_{4}^{2} = k^{2} R(t)^{2} \}
\end{equation*}
for all $t \geq 0$, and
\begin{equation*}
V^{4} = V^{4} (k, R) = \cup \{ T_{t}^{3}: 0 \leq t < \infty \}.
\end{equation*}

\begin{description}

\item[Exercise 5.15:] Show that $V^{4}$ is cross-sectionally toroidal with respect to each $t \geq 0$.

\item[Exercise 5.16:] For a fixed time $t \geq 0$, show that the intersection of $T_{t}^{3}$ with the $x_{1}x_{2}$-plane is the union of two circles of radius $(1 \pm k) R(t_{0})$. Show that if $R(t)$ is unbounded, then as $t$ varies from $0$ to $\infty$, the region swept out in the $x_{1}x_{2}$-plane is the complement of a circular disk.

\item[Exercise 5.17:] Show that the intersection of $V^{4}$ with the $x_{3}x_{4}$-plane is empty for $k < 1$.

\item[Exercise 5.18:] For $0 \leq \theta \leq 2\pi$, let
  \begin{equation*}
  P_{\theta}^{3} =
  \begin{cases}
    \{ (x_{1}, x_{1} \tan(\theta), x_{3}, 0): x_{1} \geq 0 \} \text{  if } 0 \leq \theta \frac{\pi}{2}, \text{ or } \frac{3\pi}{2} < \theta < 2\pi \\
    \{ (x_{1}, x_{1} \tan(\theta), x_{3}, 0): x_{1} \leq 0 \} \text{  if } \frac{\pi}{2} < \theta < \frac{3\pi}{2}
  \end{cases}
  \end{equation*}
  \{ cf 5.12 \}. Show that the intersection of $T_{t}^{3}$ at any time $t \geq 0$ with $P_{\theta}^{3}$, is a circle with center on the axial circle. Show that as $t$ varies from $0$ to $\infty$, the region swept out in $P_{\theta}^{3}$ is a ``triangular-shaped'' region whose boundary consists of two half-rays and part of a circle. Hint: the center of the circle is on the axial circle, and the half-rays have slope $\pm \frac{k}{\sqrt{1-k^{2}} }$.

\item[Exercise 5.19:] Define $P_{\theta}^{3}$ as in 5.18 except that $x_{4}$ is not restricted. Show that for $0<t$, the intersection $P_{\theta}^{3} \cap T_{t}^{3} = St_{\theta}^{3}$ is a $2$-sphere with center $c(t, \theta) = L_{\theta} \cap A(t)$, where $A(t)$ is the axial circle of $T_{t}^{3}$, and $L_{\theta} = P_{\theta}^{3} \cap x_{1}x_{2}$-plane is a half-ray.

\end{description}

Given a point $p \epsilon V^{4}$, there exists $t \geq 0$ such that $p = p(t) \epsilon T_{t}^{3}$, and there exists $0 \leq \theta < 2\pi$ such that $p \epsilon S_{t\theta}^{2}$ (see 5.19). The point $p(t)$ can be located by the use of four angles associated with $p$. For a given $t$, we define the \emph{angular coordinates} $(\theta, \phi, \beta, \gamma)$ of a point $p(t) \epsilon T_{t}^{3}$ as follows. The angle $\theta$ is the angle between the positive $x_{1}$-axis and the line $L_{\theta}$ (where $L_{\theta} = P_{\theta}^{3} \cap x_{1}x_{2}$-plane) measured clockwise from the positive $x_{1}$-axis. Let $L_{cp}$ be the line segment with end points $c(t, \theta)$ and $p(t)$, and let $p'$ be the perpendicular projection of $p$ onto the $x_{1}x_{2}$-plane. The angle $\phi$ is the angle between the line $L_{\theta}$ and the line $L_{cp'}$ (note that $L_{cp'}$ is the perpendicular projection of $L_{cp}$ onto the $x_{1}x_{2}$-plane) measured clockwise from $L_{\theta}$. The angle $\beta$ is defined to be the angle between the positive $x_{3}$-axis and the line segment $L_{cp}$ measured clockwise from the $x_{3}$-axis. The angle $\gamma$ is defined to be the angle between the positive $x_{4}$-axis and the line segment $L_{cp}$ measured clockwise from the $x_{4}$-axis (see Figure 3 where $\theta$, $\phi$, and $\beta$ are depicted).

%Fig3
\begin{center}
%TeXCAD Picture [meu-fig3.pic]. Options:
%\grade{\on}
%\emlines{\off}
%\epic{\off}
%\beziermacro{\on}
%\reduce{\on}
%\snapping{\off}
%\quality{8.00}
%\graddiff{0.01}
%\snapasp{1}
%\zoom{4.0000}
\unitlength 1mm % = 2.85pt
\linethickness{0.4pt}
\ifx\plotpoint\undefined\newsavebox{\plotpoint}\fi % GNUPLOT compatibility
\begin{picture}(82.5,100.5)(0,0)
%\dottedline(20.75,71)(20.75,32.5)
\multiput(20.68,70.93)(0,-.98718){40}{{\rule{.4pt}{.4pt}}}
%\end
%\dottedline(20.75,32.5)(45.25,32.5)
\multiput(20.68,32.43)(.98,0){26}{{\rule{.4pt}{.4pt}}}
%\end
%\dottedline(45.25,32.5)(45.25,71.25)
\multiput(45.18,32.43)(0,.99359){40}{{\rule{.4pt}{.4pt}}}
%\end
%\dottedline(45.25,71.25)(20.75,71.5)
\multiput(45.18,71.18)(-.98,.01){26}{{\rule{.4pt}{.4pt}}}
%\end
%\dottedline(20.75,71.5)(20.75,71.75)
\multiput(20.68,71.43)(0,.125){3}{{\rule{.4pt}{.4pt}}}
%\end
%\dottedline(20.75,71.75)(33.75,89)
\multiput(20.68,71.68)(.5652,.75){24}{{\rule{.4pt}{.4pt}}}
%\end
%\dottedline(33.75,89)(57.75,89)
\multiput(33.68,88.93)(.96,0){26}{{\rule{.4pt}{.4pt}}}
%\end
%\dottedline(57.75,89)(45.75,71.5)
\multiput(57.68,88.93)(-.5217,-.7609){24}{{\rule{.4pt}{.4pt}}}
%\end
%\dottedline(58,89)(58,88.75)
\multiput(57.93,88.93)(0,-.125){3}{{\rule{.4pt}{.4pt}}}
%\end
%\dottedline(58,88.75)(58,49.75)
\multiput(57.93,88.68)(0,-.975){41}{{\rule{.4pt}{.4pt}}}
%\end
%\dottedline(58,49.75)(45.5,32.5)
\multiput(57.93,49.68)(-.5435,-.75){24}{{\rule{.4pt}{.4pt}}}
%\end
\put(33.5,50.75){\vector(0,1){46}}
%\vector(33.75,50.75)(4.75,10.75)
\put(4.75,10.75){\vector(-3,-4){.07}}\multiput(33.75,50.75)(-.0337209302,-.0465116279){860}{\line(0,-1){.0465116279}}
%\end
\put(33.5,50.25){\vector(1,0){46.25}}
\put(34,50.25){\line(-1,0){.5}}
%\emline(33.5,50.25)(66.25,26.75)
\multiput(33.5,50.25)(.0469870875,-.0337159254){697}{\line(1,0){.0469870875}}
%\end
%\emline(38.25,46.75)(52.75,17)
\multiput(38.25,46.75)(.03372093,-.069186047){430}{\line(0,-1){.069186047}}
%\end
%\emline(38.75,47)(52.75,100.25)
\multiput(38.75,47)(.03373494,.128313253){415}{\line(0,1){.128313253}}
%\end
\put(45,71.5){\vector(0,1){29}}
%\qbezvec(45.5,92.75)(47.25,94.38)(49,90.5)
\put(49,90.5){\vector(1,-2){.07}}\qbezier(45.5,92.75)(47.25,94.38)(49,90.5)
%\end
%\qbezvec(58.75,47.75)(63.13,41.63)(54,37)
\put(54,37){\vector(-2,-1){.07}}\qbezier(58.75,47.75)(63.13,41.63)(54,37)
%\end
%\qbezvec(52.5,34)(53.88,30)(47.75,30)
\put(47.75,30){\vector(-1,0){.07}}\qbezier(52.5,34)(53.88,30)(47.75,30)
%\end
\put(62.25,41.5){\makebox(0,0)[cc]{$\theta$}}
\put(54.25,29){\makebox(0,0)[cc]{$\phi$}}
\put(42.75,28.5){\makebox(0,0)[cc]{$p'$}}
\put(5,9){\makebox(0,0)[cc]{$x_{2}$}}
\put(82.5,49.5){\makebox(0,0)[cc]{$x_{1}$}}
\put(33.5,99.25){\makebox(0,0)[cc]{$x_{3}$}}
\put(47.5,95.5){\makebox(0,0)[cc]{$\beta$}}
\put(48,70.25){\makebox(0,0)[cc]{$p$}}
\put(37,45){\makebox(0,0)[cc]{$c$}}
\put(36.25,2.75){\makebox(0,0)[cc]{Figure 3}}
\end{picture}
\end{center}

\begin{description}

\item[Exercise 5.20:] Show that the point $p(t) \epsilon T_{t}^{3}$ is determined uniquely by its angular coordinates.

\item[Exercise 5.21:] Let $p(t) \epsilon T_{t}^{3}$ have angular coordinates $(\theta, \phi, \beta, \gamma)$. For every $0 \leq s < \infty$, let $p(s)$ denote the point in $T_{s}^{3}$ with the same angular coordinates as $p(t)$. Show that
  \begin{equation*}
  \{ p(s): 0 \leq s < \infty \}
  \end{equation*}
  is a ray emanating from the origin.


\end{description}



\section{Determining Hubble's constant}
\markboth{6. Determining Hubble's constant}{6. Determining Hubble's constant}

Using the formula $H = \frac{v}{r}$ to determine Hubble's constant requires measuring $r$ and $v$, the distance to and the radial velocity of a distant galaxy, far enough away that its recession velocity is primarily caused by the expansion of the universe. Radial velocity is measured by the shift in wavelength of spectral lines from the galaxy (red shift). Various effects cause the lines to spread, resulting in uncertainties in the measurement of $v$; nevertheless, $v$ is more precisely known than $r$ (an interesting discussion of superluminal recession velocities may be found in [26, p. 242 ft]). 

Measurements of $r$ are obtained by using the inverse-square law to relate the intrinsic and apparent brightness of a galaxy. Apparent brightness is easily measured; values for intrinsic brightness are much harder to obtain (for an elementary discussion of techniques through 1976, see [24, 247-248], and more modern techniques are discussed in [23]; also see [19]). An additional complication is the necessity of modifying the inverse-square law to account for matter between the galaxy and the earth. Current estimates for $H$ lie between $50$ and $100$ (kilometers/second)/megaparsec. In the models discussed in {\S}4, it is possible to find $H$ in other ways. We show how this can be done.

Equation (7) can be solved easily for constant $H$ (note that the right-hand side of (7) is the derivative of $\log(r(t))$, and we get
\begin{equation}
r(t) = Ce^{Ht}
\end{equation}
where $C$ is a constant of integration. If we know the distance from the earth to a distant galaxy at two times (say $t_{0} < t_{1}$, and $r(t_{0}) = r_{0}$, $r(t_{1}) = r_{1})$, then we can calculate $H$ by taking the ratio $r_{1}: r_{0}$ from (8). We get
\begin{equation}
H = \frac{1}{(t_{1}-t_{0})} \ln \frac{r_{1}}{r_{0}}
\end{equation}

\begin{description}

\item[Exercise 6.1:] Assume that we know the radial velocities $v(t) = \frac{dr}{dt}(t)$ of a galaxy at two times (say $t_{0} < t_{1}$, and $v(t_{0}) = v_{0}$, $v(t_{1}) = v_{1}$). Show that
  \begin{equation}
  H = \frac{1}{(t_{1}-t_{0})} \ln \frac{v_{1}}{v_{0}}
  \end{equation}

\end{description}

We conclude this section with an order of magnitude calculation (sometimes called a Fermi calculation in physics) which shows that the formulas (9) and (10) for $H$, based on two observations of the same galaxy, cannot presently provide a more accurate value for $H$ than is already known. The cause for this lies in the relatively short time during which such observations have been made, and the limited accuracy of measurements of distances and radial velocities of galaxies. Hubble made his observations of radial velocities about $70$ years ago; so as of today, $(t_{1} - t_{0}) \approx 70$ years. Using only orders of magnitude, one megaparsec is approximately $3 \times 10^{19}$ kilometers, and $70$ years is approximately $2 \times 10^{9}$ seconds. Taking what is considered a possible ''median'' value for $H$, i.e.,
\begin{equation*}
H = 75 \text{ km/sec/Mpc } = \frac{75}{3 \times 10^{19}} \text{ sec}^{-1}
\end{equation*}
and substituting into (10), we get
\begin{equation*}
\frac{75}{3 \times 10^{19}} \approx \frac{1}{2 \times 10^{9}} \ln \frac{v_{1}}{v_{2}}
\end{equation*}
and therefore
\begin{equation*}
\frac{v_{1}}{v_{2}} \approx 1.000 000 005
\end{equation*}

Thus we would have to measure the radial velocities of galaxies to about $5$ parts per billion. Although radial velocities are more precisely known than distances, this amount of precision is not possible with present day measuring instruments. Alternately, we could wait for the time interval $(t_{1} - t_{0})$ to become large enough to compensate for the uncertainty in the radial velocities. We leave it as an exercise to determine whether this choice is practical.

\begin{description}

\item[Exercise 6.2:] Assume $\frac{v_{1}}{v_{0}} = 1.00001$ and $H = 75$. Find the number of years needed to determine this value of $H$.

\end{description}



\section{Wormholes, parallel universes, and infinite dimensions}
\markboth{7. Wormholes, Parallel Universes, etc.}{7. Wormholes, Parallel Universes, etc.}

Many science fiction stories allow spaceships to travel faster than light by ``going into hyperspace''. In the $n=2$ case, we can imagine ``hyperspace'' as a tube through the interior of the sphere. The tube connects two points on the surface, and the distance between the two points along the tube is less than the distance between the two points for paths confined to the surface of the sphere. Let us model such a tube. For brevity, we consider only the spacetime $H^{4}$ discussed above, and leave the general $n$-dimensional case to the reader. Let $p$ and $q$ be distinct points on $S_{t}^{3}$, a $3$-sphere cross-section of $H^{4}$. Let $L$ be a straight line segment with end points $p$ and $q$. Now consider an affine space $P_{L}^{4}$ that contains $L$. There are, of course, many such spaces containing $L$. Let $\delta > 0$, and define
\begin{equation*}
T(L,\delta) = \{ x \epsilon P_{L}^{4} : d(x,L) < \delta \}.
\end{equation*}

This gives the desired tube (see Figure 4). It is interesting to speculate whether matter or energy might pass through this ``hyperspace'' tube from one point of the universe $S_{t}^{3}$ to another point of $S_{t}^{3}$. The absence of an ambient space for our physical universe may make such transport impossible (however, see [17], [29], and for a less technical discussion, [6]).


%fig4
\begin{center}
%TeXCAD Picture [meu-fig4.pic]. Options:
%\grade{\on}
%\emlines{\off}
%\epic{\off}
%\beziermacro{\on}
%\reduce{\on}
%\snapping{\off}
%\quality{8.00}
%\graddiff{0.01}
%\snapasp{1}
%\zoom{4.0000}
\unitlength 1mm % = 2.85pt
\linethickness{0.4pt}
\ifx\plotpoint\undefined\newsavebox{\plotpoint}\fi % GNUPLOT compatibility
\begin{picture}(70.91,75.78)(0,0)
%\circle(37.25,42.25){67.07}
\put(70.78,42.25){\line(0,1){1.292}}
\put(70.76,43.54){\line(0,1){1.29}}
\multiput(70.69,44.83)(-.0311,.3215){4}{\line(0,1){.3215}}
\multiput(70.56,46.12)(-.02895,.21339){6}{\line(0,1){.21339}}
\multiput(70.39,47.4)(-.03184,.18181){7}{\line(0,1){.18181}}
\multiput(70.16,48.67)(-.0302,.14035){9}{\line(0,1){.14035}}
\multiput(69.89,49.93)(-.03202,.12517){10}{\line(0,1){.12517}}
\multiput(69.57,51.19)(-.03347,.11259){11}{\line(0,1){.11259}}
\multiput(69.2,52.42)(-.031972,.094105){13}{\line(0,1){.094105}}
\multiput(68.79,53.65)(-.033032,.086175){14}{\line(0,1){.086175}}
\multiput(68.33,54.85)(-.031786,.074234){16}{\line(0,1){.074234}}
\multiput(67.82,56.04)(-.032586,.068663){17}{\line(0,1){.068663}}
\multiput(67.26,57.21)(-.033251,.063614){18}{\line(0,1){.063614}}
\multiput(66.67,58.35)(-.032109,.056058){20}{\line(0,1){.056058}}
\multiput(66.02,59.48)(-.032614,.052171){21}{\line(0,1){.052171}}
\multiput(65.34,60.57)(-.033026,.048563){22}{\line(0,1){.048563}}
\multiput(64.61,61.64)(-.033356,.0452){23}{\line(0,1){.0452}}
\multiput(63.84,62.68)(-.033611,.042054){24}{\line(0,1){.042054}}
\multiput(63.04,63.69)(-.032498,.037595){26}{\line(0,1){.037595}}
\multiput(62.19,64.67)(-.032666,.03497){27}{\line(0,1){.03497}}
\multiput(61.31,65.61)(-.033989,.033686){27}{\line(-1,0){.033989}}
\multiput(60.39,66.52)(-.036617,.033596){26}{\line(-1,0){.036617}}
\multiput(59.44,67.39)(-.0394,.033447){25}{\line(-1,0){.0394}}
\multiput(58.46,68.23)(-.042353,.033233){24}{\line(-1,0){.042353}}
\multiput(57.44,69.03)(-.045497,.03295){23}{\line(-1,0){.045497}}
\multiput(56.39,69.78)(-.048857,.03259){22}{\line(-1,0){.048857}}
\multiput(55.32,70.5)(-.052461,.032145){21}{\line(-1,0){.052461}}
\multiput(54.22,71.18)(-.059309,.033269){19}{\line(-1,0){.059309}}
\multiput(53.09,71.81)(-.06391,.032679){18}{\line(-1,0){.06391}}
\multiput(51.94,72.4)(-.068952,.031969){17}{\line(-1,0){.068952}}
\multiput(50.77,72.94)(-.079483,.033195){15}{\line(-1,0){.079483}}
\multiput(49.57,73.44)(-.086467,.032259){14}{\line(-1,0){.086467}}
\multiput(48.36,73.89)(-.10225,.03372){12}{\line(-1,0){.10225}}
\multiput(47.14,74.29)(-.11288,.03246){11}{\line(-1,0){.11288}}
\multiput(45.9,74.65)(-.12546,.0309){10}{\line(-1,0){.12546}}
\multiput(44.64,74.96)(-.15819,.03255){8}{\line(-1,0){.15819}}
\multiput(43.38,75.22)(-.18209,.03021){7}{\line(-1,0){.18209}}
\multiput(42.1,75.43)(-.25636,.03245){5}{\line(-1,0){.25636}}
\put(40.82,75.59){\line(-1,0){1.287}}
\put(39.53,75.71){\line(-1,0){1.291}}
\put(38.24,75.77){\line(-1,0){1.292}}
\put(36.95,75.78){\line(-1,0){1.292}}
\put(35.66,75.75){\line(-1,0){1.289}}
\multiput(34.37,75.66)(-.25698,-.02716){5}{\line(-1,0){.25698}}
\multiput(33.08,75.53)(-.21312,-.03086){6}{\line(-1,0){.21312}}
\multiput(31.81,75.34)(-.18152,-.03347){7}{\line(-1,0){.18152}}
\multiput(30.53,75.11)(-.14007,-.03145){9}{\line(-1,0){.14007}}
\multiput(29.27,74.82)(-.12488,-.03314){10}{\line(-1,0){.12488}}
\multiput(28.03,74.49)(-.10293,-.03161){12}{\line(-1,0){.10293}}
\multiput(26.79,74.11)(-.093815,-.032814){13}{\line(-1,0){.093815}}
\multiput(25.57,73.69)(-.08015,-.031549){15}{\line(-1,0){.08015}}
\multiput(24.37,73.21)(-.073946,-.03245){16}{\line(-1,0){.073946}}
\multiput(23.19,72.69)(-.068368,-.033199){17}{\line(-1,0){.068368}}
\multiput(22.02,72.13)(-.059982,-.032039){19}{\line(-1,0){.059982}}
\multiput(20.88,71.52)(-.055768,-.03261){20}{\line(-1,0){.055768}}
\multiput(19.77,70.87)(-.051877,-.03308){21}{\line(-1,0){.051877}}
\multiput(18.68,70.17)(-.048265,-.03346){22}{\line(-1,0){.048265}}
\multiput(17.62,69.44)(-.043029,-.032353){24}{\line(-1,0){.043029}}
\multiput(16.58,68.66)(-.040081,-.032627){25}{\line(-1,0){.040081}}
\multiput(15.58,67.84)(-.037302,-.032834){26}{\line(-1,0){.037302}}
\multiput(14.61,66.99)(-.034676,-.032978){27}{\line(-1,0){.034676}}
\multiput(13.68,66.1)(-.03338,-.034289){27}{\line(0,-1){.034289}}
\multiput(12.77,65.17)(-.033266,-.036917){26}{\line(0,-1){.036917}}
\multiput(11.91,64.21)(-.033092,-.039698){25}{\line(0,-1){.039698}}
\multiput(11.08,63.22)(-.032853,-.042649){24}{\line(0,-1){.042649}}
\multiput(10.29,62.2)(-.032541,-.045791){23}{\line(0,-1){.045791}}
\multiput(9.55,61.15)(-.033682,-.051487){21}{\line(0,-1){.051487}}
\multiput(8.84,60.06)(-.033257,-.055384){20}{\line(0,-1){.055384}}
\multiput(8.17,58.96)(-.032736,-.059604){19}{\line(0,-1){.059604}}
\multiput(7.55,57.82)(-.032105,-.0642){18}{\line(0,-1){.0642}}
\multiput(6.97,56.67)(-.03331,-.073563){16}{\line(0,-1){.073563}}
\multiput(6.44,55.49)(-.032481,-.079777){15}{\line(0,-1){.079777}}
\multiput(5.95,54.3)(-.031483,-.086753){14}{\line(0,-1){.086753}}
\multiput(5.51,53.08)(-.0328,-.10255){12}{\line(0,-1){.10255}}
\multiput(5.12,51.85)(-.03145,-.11317){11}{\line(0,-1){.11317}}
\multiput(4.77,50.61)(-.03308,-.1397){9}{\line(0,-1){.1397}}
\multiput(4.47,49.35)(-.03114,-.15848){8}{\line(0,-1){.15848}}
\multiput(4.23,48.08)(-.03334,-.21274){6}{\line(0,-1){.21274}}
\multiput(4.03,46.8)(-.03015,-.25664){5}{\line(0,-1){.25664}}
\put(3.87,45.52){\line(0,-1){1.288}}
\put(3.77,44.23){\line(0,-1){5.163}}
\multiput(3.87,39.07)(.02946,-.25672){5}{\line(0,-1){.25672}}
\multiput(4.01,37.79)(.03277,-.21283){6}{\line(0,-1){.21283}}
\multiput(4.21,36.51)(.03071,-.15856){8}{\line(0,-1){.15856}}
\multiput(4.46,35.24)(.03271,-.13979){9}{\line(0,-1){.13979}}
\multiput(4.75,33.98)(.03114,-.11325){11}{\line(0,-1){.11325}}
\multiput(5.09,32.74)(.03253,-.10264){12}{\line(0,-1){.10264}}
\multiput(5.48,31.51)(.033653,-.093517){13}{\line(0,-1){.093517}}
\multiput(5.92,30.29)(.032266,-.079865){15}{\line(0,-1){.079865}}
\multiput(6.4,29.09)(.033111,-.073652){16}{\line(0,-1){.073652}}
\multiput(6.93,27.91)(.031932,-.064286){18}{\line(0,-1){.064286}}
\multiput(7.51,26.76)(.032575,-.059692){19}{\line(0,-1){.059692}}
\multiput(8.13,25.62)(.033108,-.055474){20}{\line(0,-1){.055474}}
\multiput(8.79,24.51)(.033543,-.051578){21}{\line(0,-1){.051578}}
\multiput(9.49,23.43)(.032418,-.045878){23}{\line(0,-1){.045878}}
\multiput(10.24,22.37)(.032737,-.042737){24}{\line(0,-1){.042737}}
\multiput(11.03,21.35)(.032985,-.039787){25}{\line(0,-1){.039787}}
\multiput(11.85,20.35)(.033167,-.037006){26}{\line(0,-1){.037006}}
\multiput(12.71,19.39)(.033287,-.034379){27}{\line(0,-1){.034379}}
\multiput(13.61,18.46)(.034587,-.033071){27}{\line(1,0){.034587}}
\multiput(14.55,17.57)(.037213,-.032934){26}{\line(1,0){.037213}}
\multiput(15.51,16.71)(.039993,-.032735){25}{\line(1,0){.039993}}
\multiput(16.51,15.9)(.042942,-.032469){24}{\line(1,0){.042942}}
\multiput(17.54,15.12)(.048175,-.03359){22}{\line(1,0){.048175}}
\multiput(18.6,14.38)(.051787,-.033219){21}{\line(1,0){.051787}}
\multiput(19.69,13.68)(.05568,-.03276){20}{\line(1,0){.05568}}
\multiput(20.8,13.02)(.059895,-.032201){19}{\line(1,0){.059895}}
\multiput(21.94,12.41)(.068278,-.033384){17}{\line(1,0){.068278}}
\multiput(23.1,11.85)(.073858,-.032649){16}{\line(1,0){.073858}}
\multiput(24.28,11.32)(.080065,-.031765){15}{\line(1,0){.080065}}
\multiput(25.49,10.85)(.093726,-.033066){13}{\line(1,0){.093726}}
\multiput(26.7,10.42)(.10284,-.03188){12}{\line(1,0){.10284}}
\multiput(27.94,10.03)(.12479,-.03348){10}{\line(1,0){.12479}}
\multiput(29.19,9.7)(.13999,-.03183){9}{\line(1,0){.13999}}
\multiput(30.45,9.41)(.15875,-.02972){8}{\line(1,0){.15875}}
\multiput(31.72,9.17)(.21303,-.03144){6}{\line(1,0){.21303}}
\multiput(32.99,8.99)(.2569,-.02785){5}{\line(1,0){.2569}}
\put(34.28,8.85){\line(1,0){1.289}}
\put(35.57,8.76){\line(1,0){1.291}}
\put(36.86,8.72){\line(1,0){1.292}}
\put(38.15,8.73){\line(1,0){1.291}}
\put(39.44,8.79){\line(1,0){1.287}}
\multiput(40.73,8.9)(.25645,.03176){5}{\line(1,0){.25645}}
\multiput(42.01,9.05)(.18217,.02972){7}{\line(1,0){.18217}}
\multiput(43.29,9.26)(.15828,.03213){8}{\line(1,0){.15828}}
\multiput(44.55,9.52)(.12554,.03056){10}{\line(1,0){.12554}}
\multiput(45.81,9.83)(.11297,.03216){11}{\line(1,0){.11297}}
\multiput(47.05,10.18)(.10234,.03345){12}{\line(1,0){.10234}}
\multiput(48.28,10.58)(.086554,.032026){14}{\line(1,0){.086554}}
\multiput(49.49,11.03)(.079572,.03298){15}{\line(1,0){.079572}}
\multiput(50.68,11.52)(.069038,.031783){17}{\line(1,0){.069038}}
\multiput(51.86,12.06)(.063998,.032507){18}{\line(1,0){.063998}}
\multiput(53.01,12.65)(.059398,.033109){19}{\line(1,0){.059398}}
\multiput(54.14,13.28)(.055175,.033603){20}{\line(1,0){.055175}}
\multiput(55.24,13.95)(.048945,.032458){22}{\line(1,0){.048945}}
\multiput(56.32,14.66)(.045586,.032827){23}{\line(1,0){.045586}}
\multiput(57.37,15.42)(.042442,.033119){24}{\line(1,0){.042442}}
\multiput(58.39,16.21)(.03949,.03334){25}{\line(1,0){.03949}}
\multiput(59.37,17.05)(.036708,.033497){26}{\line(1,0){.036708}}
\multiput(60.33,17.92)(.034079,.033594){27}{\line(1,0){.034079}}
\multiput(61.25,18.83)(.03276,.034882){27}{\line(0,1){.034882}}
\multiput(62.13,19.77)(.0326,.037507){26}{\line(0,1){.037507}}
\multiput(62.98,20.74)(.033725,.041963){24}{\line(0,1){.041963}}
\multiput(63.79,21.75)(.033478,.04511){23}{\line(0,1){.04511}}
\multiput(64.56,22.79)(.033157,.048474){22}{\line(0,1){.048474}}
\multiput(65.29,23.85)(.032754,.052083){21}{\line(0,1){.052083}}
\multiput(65.98,24.95)(.03226,.055971){20}{\line(0,1){.055971}}
\multiput(66.62,26.07)(.033422,.063525){18}{\line(0,1){.063525}}
\multiput(67.22,27.21)(.032771,.068575){17}{\line(0,1){.068575}}
\multiput(67.78,28.38)(.031986,.074148){16}{\line(0,1){.074148}}
\multiput(68.29,29.56)(.033264,.086085){14}{\line(0,1){.086085}}
\multiput(68.76,30.77)(.032225,.094019){13}{\line(0,1){.094019}}
\multiput(69.18,31.99)(.03096,.10312){12}{\line(0,1){.10312}}
\multiput(69.55,33.23)(.03236,.12509){10}{\line(0,1){.12509}}
\multiput(69.87,34.48)(.03057,.14027){9}{\line(0,1){.14027}}
\multiput(70.15,35.74)(.03233,.18172){7}{\line(0,1){.18172}}
\multiput(70.37,37.01)(.02953,.21331){6}{\line(0,1){.21331}}
\multiput(70.55,38.29)(.0319,.3214){4}{\line(0,1){.3214}}
\put(70.68,39.58){\line(0,1){1.29}}
\put(70.76,40.87){\line(0,1){1.382}}
%\end
\qbezier(3.5,42.25)(38.88,23.38)(70.75,42)
\put(36.5,2.5){\makebox(0,0)[cc]{Figure 4}}
\qbezier(24.68,61.54)(24.75,57.08)(19.77,58.57)
\qbezier(19.62,58.72)(20.07,64.14)(24.68,61.54)
%\dashline{1}(19.62,58.57)(45.64,27.2)
\multiput(19.55,58.5)(.033618,-.04053){18}{\line(0,-1){.04053}}
\multiput(20.76,57.04)(.033618,-.04053){18}{\line(0,-1){.04053}}
\multiput(21.97,55.58)(.033618,-.04053){18}{\line(0,-1){.04053}}
\multiput(23.18,54.12)(.033618,-.04053){18}{\line(0,-1){.04053}}
\multiput(24.39,52.66)(.033618,-.04053){18}{\line(0,-1){.04053}}
\multiput(25.6,51.2)(.033618,-.04053){18}{\line(0,-1){.04053}}
\multiput(26.81,49.75)(.033618,-.04053){18}{\line(0,-1){.04053}}
\multiput(28.02,48.29)(.033618,-.04053){18}{\line(0,-1){.04053}}
\multiput(29.23,46.83)(.033618,-.04053){18}{\line(0,-1){.04053}}
\multiput(30.44,45.37)(.033618,-.04053){18}{\line(0,-1){.04053}}
\multiput(31.65,43.91)(.033618,-.04053){18}{\line(0,-1){.04053}}
\multiput(32.86,42.45)(.033618,-.04053){18}{\line(0,-1){.04053}}
\multiput(34.07,40.99)(.033618,-.04053){18}{\line(0,-1){.04053}}
\multiput(35.28,39.53)(.033618,-.04053){18}{\line(0,-1){.04053}}
\multiput(36.49,38.07)(.033618,-.04053){18}{\line(0,-1){.04053}}
\multiput(37.7,36.61)(.033618,-.04053){18}{\line(0,-1){.04053}}
\multiput(38.91,35.15)(.033618,-.04053){18}{\line(0,-1){.04053}}
\multiput(40.12,33.7)(.033618,-.04053){18}{\line(0,-1){.04053}}
\multiput(41.33,32.24)(.033618,-.04053){18}{\line(0,-1){.04053}}
\multiput(42.54,30.78)(.033618,-.04053){18}{\line(0,-1){.04053}}
\multiput(43.75,29.32)(.033618,-.04053){18}{\line(0,-1){.04053}}
\multiput(44.96,27.86)(.033618,-.04053){18}{\line(0,-1){.04053}}
%\end
%\dashline{1}(24.68,61.54)(50.69,30.32)
\multiput(24.61,61.47)(.033605,-.040336){18}{\line(0,-1){.040336}}
\multiput(25.82,60.02)(.033605,-.040336){18}{\line(0,-1){.040336}}
\multiput(27.03,58.57)(.033605,-.040336){18}{\line(0,-1){.040336}}
\multiput(28.24,57.11)(.033605,-.040336){18}{\line(0,-1){.040336}}
\multiput(29.45,55.66)(.033605,-.040336){18}{\line(0,-1){.040336}}
\multiput(30.66,54.21)(.033605,-.040336){18}{\line(0,-1){.040336}}
\multiput(31.87,52.76)(.033605,-.040336){18}{\line(0,-1){.040336}}
\multiput(33.08,51.31)(.033605,-.040336){18}{\line(0,-1){.040336}}
\multiput(34.29,49.85)(.033605,-.040336){18}{\line(0,-1){.040336}}
\multiput(35.5,48.4)(.033605,-.040336){18}{\line(0,-1){.040336}}
\multiput(36.71,46.95)(.033605,-.040336){18}{\line(0,-1){.040336}}
\multiput(37.92,45.5)(.033605,-.040336){18}{\line(0,-1){.040336}}
\multiput(39.13,44.04)(.033605,-.040336){18}{\line(0,-1){.040336}}
\multiput(40.34,42.59)(.033605,-.040336){18}{\line(0,-1){.040336}}
\multiput(41.55,41.14)(.033605,-.040336){18}{\line(0,-1){.040336}}
\multiput(42.76,39.69)(.033605,-.040336){18}{\line(0,-1){.040336}}
\multiput(43.97,38.24)(.033605,-.040336){18}{\line(0,-1){.040336}}
\multiput(45.18,36.78)(.033605,-.040336){18}{\line(0,-1){.040336}}
\multiput(46.39,35.33)(.033605,-.040336){18}{\line(0,-1){.040336}}
\multiput(47.6,33.88)(.033605,-.040336){18}{\line(0,-1){.040336}}
\multiput(48.81,32.43)(.033605,-.040336){18}{\line(0,-1){.040336}}
\multiput(50.01,30.98)(.033605,-.040336){18}{\line(0,-1){.040336}}
%\end
%\dashline{1}(3.86,42.51)(7.58,44.74)
\multiput(3.79,42.44)(.053089,.031854){14}{\line(1,0){.053089}}
\multiput(5.28,43.34)(.053089,.031854){14}{\line(1,0){.053089}}
\multiput(6.77,44.23)(.053089,.031854){14}{\line(1,0){.053089}}
%\end
%\dashline{1}(7.58,44.74)(12.04,46.53)
\multiput(7.51,44.67)(.08258,.03303){9}{\line(1,0){.08258}}
\multiput(9,45.27)(.08258,.03303){9}{\line(1,0){.08258}}
\multiput(10.48,45.86)(.08258,.03303){9}{\line(1,0){.08258}}
%\end
%\dashline{1}(12.04,46.53)(18.28,48.61)
\multiput(11.97,46.46)(.09755,.03252){8}{\line(1,0){.09755}}
\multiput(13.53,46.98)(.09755,.03252){8}{\line(1,0){.09755}}
\multiput(15.09,47.5)(.09755,.03252){8}{\line(1,0){.09755}}
\multiput(16.65,48.02)(.09755,.03252){8}{\line(1,0){.09755}}
%\end
%\dashline{1}(18.28,48.61)(24.08,49.95)
\multiput(18.21,48.54)(.13803,.03185){6}{\line(1,0){.13803}}
\multiput(19.87,48.92)(.13803,.03185){6}{\line(1,0){.13803}}
\multiput(21.53,49.3)(.13803,.03185){6}{\line(1,0){.13803}}
\multiput(23.18,49.69)(.13803,.03185){6}{\line(1,0){.13803}}
%\end
%\dashline{1}(24.08,49.95)(26.46,50.24)
\put(24.01,49.88){\line(1,0){.793}}
\put(25.6,50.07){\line(1,0){.793}}
%\end
%\dashline{1}(33.3,50.99)(38.35,50.99)
\put(33.23,50.92){\line(1,0){.842}}
\put(34.91,50.92){\line(1,0){.842}}
\put(36.6,50.92){\line(1,0){.842}}
%\end
%\dashline{1}(38.35,50.99)(49.65,49.8)
\put(38.28,50.92){\line(1,0){.941}}
\put(40.16,50.72){\line(1,0){.941}}
\put(42.05,50.52){\line(1,0){.941}}
\put(43.93,50.32){\line(1,0){.941}}
\put(45.81,50.12){\line(1,0){.941}}
\put(47.7,49.93){\line(1,0){.941}}
%\end
%\dashline{1}(49.65,49.8)(55.6,48.16)
\multiput(49.58,49.73)(.12135,-.03337){7}{\line(1,0){.12135}}
\multiput(51.28,49.26)(.12135,-.03337){7}{\line(1,0){.12135}}
\multiput(52.98,48.79)(.12135,-.03337){7}{\line(1,0){.12135}}
\multiput(54.68,48.33)(.12135,-.03337){7}{\line(1,0){.12135}}
%\end
%\dashline{1}(55.6,48.16)(62.58,45.78)
\multiput(55.52,48.09)(.09704,-.03303){9}{\line(1,0){.09704}}
\multiput(57.27,47.5)(.09704,-.03303){9}{\line(1,0){.09704}}
\multiput(59.02,46.9)(.09704,-.03303){9}{\line(1,0){.09704}}
\multiput(60.76,46.31)(.09704,-.03303){9}{\line(1,0){.09704}}
%\end
%\dashline{1}(62.58,45.78)(65.7,44.74)
\multiput(62.51,45.71)(.08919,-.02973){7}{\line(1,0){.08919}}
\multiput(63.76,45.3)(.08919,-.02973){7}{\line(1,0){.08919}}
\multiput(65.01,44.88)(.08919,-.02973){7}{\line(1,0){.08919}}
%\end
%\dashline{1}(65.7,44.74)(70.16,42.66)
\multiput(65.63,44.67)(.06757,-.03153){11}{\line(1,0){.06757}}
\multiput(67.12,43.98)(.06757,-.03153){11}{\line(1,0){.06757}}
\multiput(68.61,43.29)(.06757,-.03153){11}{\line(1,0){.06757}}
%\end
%\dashline{1}(70.16,42.66)(70.91,42.22)
\multiput(70.09,42.59)(.05309,-.03185){7}{\line(1,0){.05309}}
%\end
%\dashline{1}(45.49,26.76)(45.34,27.65)
\put(45.42,26.69){\line(0,1){.446}}
%\end
%\dashline{1}(45.34,27.65)(45.49,29.14)
\put(45.27,27.58){\line(0,1){.743}}
%\end
%\dashline{1}(45.49,29.14)(45.93,30.03)
\multiput(45.42,29.07)(.03185,.06371){7}{\line(0,1){.06371}}
%\end
%\dashline{1}(45.93,30.03)(47.72,30.92)
\multiput(45.86,29.96)(.063707,.031854){14}{\line(1,0){.063707}}
%\end
%\dashline{1}(47.72,30.92)(49.5,30.62)
\multiput(47.65,30.85)(.17838,-.02973){5}{\line(1,0){.17838}}
%\end
%\dashline{1}(49.5,30.62)(50.69,30.03)
\multiput(49.43,30.55)(.06607,-.03303){9}{\line(1,0){.06607}}
%\end
%\dashline{1}(50.69,30.03)(50.1,27.8)
\multiput(50.62,29.96)(-.03303,-.12388){6}{\line(0,-1){.12388}}
\multiput(50.22,28.47)(-.03303,-.12388){6}{\line(0,-1){.12388}}
%\end
%\dashline{1}(50.1,27.8)(48.46,26.76)
\multiput(50.02,27.73)(-.04955,-.03153){11}{\line(-1,0){.04955}}
\multiput(48.93,27.03)(-.04955,-.03153){11}{\line(-1,0){.04955}}
%\end
%\dashline{1}(48.46,26.76)(47.12,26.76)
\put(48.39,26.69){\line(-1,0){.669}}
%\end
%\dashline{1}(47.12,26.76)(45.34,26.91)
\put(47.05,26.69){\line(-1,0){.892}}
%\end
\end{picture}
\end{center}

In studying the foundations of quantum mechanics, physicists have been led to a number of results that contravene common sense. One famous such paradox is the experiment which involves Schr�inger's cat. The cat is put into a sealed box with a cyanide capsule. If a radioactive atom decays, the capsule will break, and the cat will die. Because the decay of a single atom cannot be predicted, the outcome of the experiment is not known until the box is opened. Using the mathematics of quantum mechanics, the situation inside the box before it is opened is described as a superposition of two wave functions. One describes a dead cat, and the other a live cat. When we open the box, the superposition ``collapses'' to one or the other function, depending on the result we see.

Common sense tells us that cat cannot be partly alive and partly dead before we open the box. It must be one or the other, and hence the supposed wave function cannot describe ``reality''. There are various ways of reconciling quantum mechanics and common sense. One of the more exotic ways declares that when the experiment is done, two parallel universes are created. In one the cat dies, and in the other it lives (see [2] and [20, Figure 20]). For another way in which parallel universes might arise, see [28].

We now construct a model in which there are infinitely many spacetimes (``parallel universes''). While this can be done in finite dimensions, we take this opportunity to introduce an infinite dimensional version of Euclidean space which is a subset of the well-known Banach space $l^{2}$.

The space $l^{2}$ is the set of all countable-tuples (also called $\omega$-tuples)
\begin{equation*}
x = (x_{1}, x_{2}, \cdots, x_{n}, \cdots)
\end{equation*}
of real numbers with the property that the infinite series $\Sigma x_{i}^{2}$ converges. We define $E^{\infty}$ to be the subset of $l^{2}$ consisting of all points $x$ in $l^{2}$ for which there exists an integer $n$ such that $x_{i} = 0$ for all $i>n$ (the notation $E^{\infty}$ follows the convention we have been using, that superscripts indicate dimension). The point with all coordinates equal to $0$ is called the origin in $E^{\infty}$.

Each Euclidean space $E^{n}$ can be naturally embedded into $E^{\infty}$ in many ways. For example, the mapping defined by
\begin{equation}
h_{0}^{n} (x_{1}, x_{2}, \cdots, x_{n}) = (x_{1}, x_{2}, \cdots, x_{n}, 0, 0, \cdots)
\end{equation}
embeds $E^{n}$ onto
\begin{equation}
\{ x \epsilon \Sigma^{\infty} : x_{i} = 0 \text{ for } i > n \}
\end{equation}

It will be clear from our discussion that in $E^{\infty}$, we could easily allow different spacetimes to have different dimensions, and even have different time axes, but for the sake of brevity, all the spacetimes will be copies of $H^{4}$ discussed in {\S}4 and share the same time axis.

In order to get infinitely many copies of $H^{4} \subset E^{5}$, we partition the set $\{ 2, 3, 4, \cdots, n, \cdots \}$ into consecutive sets $J_{j}$ of size $4$ where

\begin{center}
$J_{j}$ = {\{}4$j$ + 2, 4$j$ + 3, 4$j$ + 4, 4$j$ + 5{\}} (0 $\leq $ j).
\end{center}

Thus $J_{o}$ = {\{}2, 3, 4, 5{\}}, $J_{1}$ = {\{}6, 7, 8, 9{\}},
$J_{2}$ = {\{}10, 11, 12, 13{\}} and so on. For each $j$ $\geq $ 0, define

\begin{center}
$X_{j}^{5}$ = {\{}$x$ $\epsilon $ $E^{\infty }$ : $x_{i}$ = 0
unless $i$ = 1, or $i$ $\epsilon$ $J_{j}${\}}.
\end{center}

For example, a point $(x_{1}, x_{2}, \cdots, x_{n}, \cdots)$ in $X_{2}^{5}$ has all $x_{i} = 0$, except possibly $x_{10}$, $x_{11}$, $x_{12}$, $x_{13}$, and $x_{1}$. To see that each $X_{j}^{5}$ is a copy of $E^{5}$, we use the following mappings. For $j \geq 0$, define $h_{j}^{5} : E^{5} \rightarrow X_{j}^{5}$ by

\begin{equation*}
h_{j}^{5} (x_{1}, x_{2}, x_{3}, x_{4}, x_{5}) = (x_{1}, 0, \cdots, 0, x_{2}, x_{3}, x_{4}, x_{5}, 0, 0, \cdots).
\end{equation*}

\begin{center}
%TeXCAD Picture [meu-symbol.pic]. Options:
%\grade{\on}
%\emlines{\off}
%\epic{\off}
%\beziermacro{\on}
%\reduce{\on}
%\snapping{\off}
%\quality{8.00}
%\graddiff{0.01}
%\snapasp{1}
%\zoom{4.0000}
\unitlength 1mm % = 2.85pt
\linethickness{0.4pt}
\ifx\plotpoint\undefined\newsavebox{\plotpoint}\fi % GNUPLOT compatibility
\begin{picture}(15.32,9.42)(81,144)
%\emline(82.57,153.42)(89.57,152.17)
\multiput(82.57,153.42)(.1842105,-.0328947){38}{\line(1,0){.1842105}}
%\end
%\emline(89.57,152.17)(96.32,153.42)
\multiput(89.57,152.17)(.1776316,.0328947){38}{\line(1,0){.1776316}}
%\end
\put(90.07,148.17){\makebox(0,0)[cc]{$4j+1$}}
\end{picture}
\end{center}

Thus, $h_{0}^{5}$ is the map defined in (11), and $X_{0}^{5}$ the set defined in (12) with $n=5$. Further, $h_{1}^{5}$ maps $E^{5}$ onto $X_{1}^{5}$, $h_{2}^{5}$ maps $E^{5}$ onto $X_{2}^{5}$, and so on. The sets $X_{j}^{5}$ have only the $x_{1}$-axis of $E^{\infty}$ in common (the time axis). For each $j$, let $H_{j}^{4}$ be the images of $H^{4}$ under the mapping $h_{j}^{5}$.

Thus $H_{j}^{4}$ is a copy of the manifold $H^{4}$ in $X_{j}^{5}$, and $X_{j}^{5}$ serves as the ambient space of $H_{j}^{4}$. Moreover, $E^{\infty}$ serves as an ambient space for the set of all ambient spaces $X_{j}^{5}$ ($j = 0, 1, 2, \cdots$). The family of $H_{j}^{4}$ ($j = 1, 2, \cdots$) gives the desired model.

It is natural to ask whether we can connect these universes by ``tubes,'' like the ``hyperspace'' tubes between points within a universe that we constructed earlier. The answer is yes; there are many ways that tubes can be constructed, and we give a simple example. For two distinct points $p, q$ with $p$ in $H_{j}^{4}$ and $q$ in $H_{4}^{k}$ ($k \neq j$), let
\begin{equation*}
X_{j,k}^{9} = \{ x \epsilon E^{\infty} : x_{i} = 0 \text{ unless } i \epsilon \{ 1 \} \cup J_{j} \cup J_{k} \}
\end{equation*}
and let $L$ be a straight line segment in $X_{j,k}^{9}$ with end points $p$ and $q$, i.e.,
\begin{equation*}
L = \{ \lambda + (1 - \lambda) q : 0 \leq \lambda \leq 1 \}.
\end{equation*}

Let $P_{L}^{4}$ be an affine space in $X_{j,k}^{9}$ that contains $L$. The affine space $P_{L}^{4}$ is a copy of $E^{4}$ in $X_{j,k}^{9}$ that may or may not contain the origin of $X_{j,k}^{9}$. Let $\delta > 0$, and define
\begin{equation*}
T(L,\delta) = \{ x \epsilon P_{L}^{4}: d(x, L) < \delta \}.
\end{equation*}

We call $T(L, \delta)$ a \emph{tube} in $E^{\infty}$ from the spacetime $H_{j}^{4}$ to the spacetime $H_{k}^{4}$. If $m$ is neither of $k$ or $j$, then obviously $T(L, \delta)$ and $X_{m}^{5}$ have at most the origin in common, and therefore the tube $T(L, \delta)$ connects two of the parallel spacetime manifolds without going through any of the other spacetime manifolds. We leave it to the reader to define tubes with other properties.

The transport of matter and energy through ``hyperspace'' tubes within a universe and tubes between universes is discussed in [17], [29], and [6] (see [22, Chapter 8] for a popularization). Both types of tubes, referred to as ``wormholes'' in physics, require the existence of matter with negative energy density. No material with this property is known to exist, but the existence and implications of such material are currently being considered in physics. The quantum phenomena permitting the possibility of such matter are discussed in [17] and [6]. For another implication of such matter, see [27].

\begin{description}

\item[Exercise 7.1:] Describe a model in which for each positive integer $n$, there is a spacetime of dimension $n$.

\item[Exercise 7.2:] Describe a model in which each of infinitely many spacetimes have different time axes.

\item[Exercise 7.3:] Describe a finite-dimensional model in which there are infinitely many parallel universes.


\end{description}



\section*{References}
\addcontentsline{toc}{section}{\numberline{}References}
\markboth{References}{References}

\begin{enumerate}[1.]
  \item Edwin A. Abbot, \emph{Flatland}, Barnes and Noble Inc., New York, 1963.
  \item Yoav Ben-Dov, \emph{Everett's theory and the ``many-worlds'' interpretation}, Am. J. Phys. 58 (9) (1990) p829-832.
  \item Donald W. Blackett, \emph{Elementary Topology}, Academic Press, New York, 1967.
  \item Dionys Burger, \emph{Sphereland}, Thomas Y. Crowell Co., New York, 1965.
  \item Ronald W. Clark, \emph{Einstein, The Life and Times}, The World Pub. Co., New York, 1971.
  \item John G. Cramer, \emph{Natural wormholes: squeezing the vacuum}, Analog Science Fiction and Fact, (July, 1992) p182-187.
  \item John G. Cramer, \emph{Stretch marks of the universe}, Analog Science Fiction and Fact, (November, 1994) p110-113.
  \item Wendy L. Freedman, \emph{The expansion rate and size of the universe}, Scientific American, (November, 1992) p54-60.
  \item Daniel. Z. Freeman and Peter van Nieuwenhuizen, \emph{The hidden dimensions of spacetime}, Scientific American, 25, (March, 1985) p74-81.
  \item George Gamow, \emph{My World Line}, Viking Press, New York, 1970.
  \item Michael B. Green, John H. Schwarz, and Edward Witten, \emph{Superstring Theory}, Cambridge University Press, Cambridge, 1987.
  \item Alex Harvey, \emph{Cosmological models}, Am. J. Phys. 61 (10) (1993), p901-906.
  \item S. W. Hawking and G. F. R. Ellis, \emph{The Large Scale Structure of Spacetime}, Cambridge University Press, Cambridge, England, 1973.
  \item E. Hubble, \emph{A relation between distance and radial velocity among extra-galactic nebulae}, Proc. Nat. Acad. Sci., 15 (1926) p169-173.
  \item D. H. Menzel, F. L. Whipple, and G. de Vaucouleurs, \emph{Survey of the Universe}, Prentice Hall, 1970.
  \item C. W. Misner, K. S. Thorne, J. A. Wheeler, \emph{Gravitation}, W.H. Freeman and Co., San Francisco, 1973.
  \item M. S. Morris and K. S. Thorne, \emph{Wormholes in spacetime and their use for interstellar travel: A tool for teaching general relativity}, Am. J. Physics 56 (5) (1988), p395-412.
  \item Gregory L. Naber, \emph{Spacetime and Singularities, An Introduction}, London Mathematical Society Student Text 11, Cambridge University Press, Cambridge, England, 1988.
  \item Donald E. Osterbrock, Joel A. Gwinn, and Ronald S. Brashea. \emph{Edwin Hubble and the expanding universe}, Scientific American, (July 1993) p84-90.
  \item Paul Renteln, \emph{Quantum Gravity}, American Scientist 79, (Nov-Dec, 1991) p508-527.
  \item Rudolf v. B. Rucker, \emph{Geometry, Relativity, and the Fourth Dimension}, Dover Publications Inc., New York, 1977.
  \item Rudolf v. B. Rucker, \emph{The Fourth Dimension}, Houghton Mifflin Co., Boston, 1984.
  \item Bertram Schwarzschild, \emph{Supernova distance measurements suggest an older, larger universe}, Physics Today, (November, 1992) p17-20.
  \item H. L. Shipman, \emph{Black Holes, Quasars and the Universe}, Houghton-Mifflin Co., Boston, 1976.
  \item R. Smith, \emph{The Expanding Universe}, Cambridge University Press, Cambridge, England, 1982.
  \item W. M. Stuckey, \emph{Can galaxies exist within our particle horizon with Hubble recessional velocities greater than $c$?} Am. J. Phys. 60 (1992) p142-146.
  \item Micllael Szpir, \emph{Spacetime Hypersurfing?}, American Scientist (82) (September-October, 1994), p422-423.
  \item Shana Tribiano, \emph{Universes without end}, The Sciences, (October 1994), p10-12.
  \item A. Zichichi, V. de Sabbata, and N. Sanches, editors, \emph{Gravitation and Modern Cosmology}, Plenum Press, New York, 1991.
\end{enumerate}

\end{document}
